\documentclass[11pt]{article}

\usepackage[T1]{fontenc}
\usepackage[utf8]{inputenc}
\usepackage[margin=1in]{geometry}
\usepackage{microtype}
\usepackage{amsmath}
\usepackage{amssymb}
\usepackage{amsthm}
\usepackage{mathtools}
\usepackage{bm}
\usepackage{booktabs}
\usepackage{array}
\usepackage{enumitem}
\usepackage[numbers,sort&compress]{natbib}
\usepackage[hidelinks]{hyperref}

\setlength{\parindent}{0pt}
\setlength{\parskip}{0.65em}
\setcounter{tocdepth}{2}

\newtheorem{theorem}{Theorem}[section]
\newtheorem{proposition}[theorem]{Proposition}
\newtheorem{lemma}[theorem]{Lemma}
\newtheorem{corollary}[theorem]{Corollary}
\newtheorem{conjecture}[theorem]{Conjecture}
\newtheorem{definition}[theorem]{Definition}
\newtheorem{remark}[theorem]{Remark}

\newcommand{\Gfifteen}{G_{15}}
\newcommand{\Gthirty}{G_{30}}
\newcommand{\Gsixty}{G_{60}}
\newcommand{\Petersen}{P_{10}}
\newcommand{\one}{\mathbf{1}}
\newcommand{\Pzero}{P_{0}}
\newcommand{\Pperp}{P_{\perp}}
\newcommand{\Rrel}{R}
\newcommand{\Ccommon}{C}
\newcommand{\zetaRel}{\zeta}
\newcommand{\tauT}{\tau_{\mathrm{T}}}
\newcommand{\spec}{\operatorname{spec}}

\title{
Contrast Concentration and the\\
Thalean Relational Saturation Conjecture\\[0.4em]
\large A Finite Linear Theorem on
$L(\mathrm{Petersen})$
}

\author{
Scott Allen Cave\\
Center of Recursive Inquiry (CoRI)\\
Matsqui Territory -- Abbotsford, British Columbia, Canada\\
\texttt{scott@rook.ca}\\
ORCID: 0009-0000-0888-4232
}

\date{August 2026\\[0.3em]
\small Reserved DOI: \href{https://doi.org/10.5281/zenodo.22210907}{10.5281/zenodo.22210907}}

\begin{document}

\maketitle

\begin{abstract}
We isolate an exact finite linear mechanism beneath the Thalean
Relational Saturation Conjecture.  On the quotient core
\[
\Gfifteen \cong L(\Petersen),
\]
the transport-induced sector--edge incidence matrix
\(M\in\{0,1\}^{15\times 30}\) defines the Thalean quadratic
\[
Q:=MM^{\mathsf T}=A^{3}+2A^{2}+2I,
\]
where \(A\) is the adjacency matrix of \(\Gfifteen\).
The common-mode line has \(Q\)-eigenvalue \(98\), while every
centered eigenspace has eigenvalue at most \(18\).
We prove that repeated algebraic aggregation by \(Q\) causes
centered relational contrast to become exponentially small relative
to the common aggregate, with contraction ratio at most
\[
\frac{18}{98}=\frac{9}{49}
\]
per iterate.  At the same time, \(Q\) is invertible, so complete
linear distinctions are not merged by this aggregation.

This establishes an exact finite instance of strong aggregate
growth together with weakening relative individuation.  We then
formulate projection-relative relational saturation: a regime in
which complete distinctions or retained histories remain nontrivial
while a declared visible or scalar readout becomes increasingly
common-mode dominated.  The conjectural extension requires an
additional native crosswalk showing that an admitted Thalean loading
sequence realizes \(Q\)-iteration, a normalized \(Q\)-type operator,
or another certified contrast-concentrating law.  Local
admissibility is kept distinct from global compatibility, and
failure to construct a globally closed inter-face transport relation
is not interpreted as saturation.

The result is entirely finite and internal.  No identification is
made with physical gravity, black holes, spacetime horizons,
singularity avoidance, Planck-scale information density, or the
cosmic microwave background.
\end{abstract}

\tableofcontents

\section{Introduction}

A finite linear system can preserve every exact state distinction and
nevertheless become progressively less informative under a declared
readout.

That elementary observation becomes structurally nontrivial in the
finite relational mechanics built on
\[
G_{60}\cong \mathrm{AT4val}[60,6].
\]
The established construction already distinguishes quotient-visible
state, lifted state, covering receipt, registered history, additive
winding, and centered scalar magnitude. Projection repeatedly forgets
finite information while more resolved constructions retain distinctions
that are invisible after quotienting or scalar compression.

The present paper asks a narrower question.

Can the exact incidence algebra on
\[
G_{15}\cong L(\mathrm{Petersen})
\]
produce a mathematically controlled concentration of relational
contrast while complete linear distinction remains preserved?

The answer is yes.

Let
\[
M\in\{0,1\}^{15\times30}
\]
be the transport-induced sector-edge incidence matrix and define
\[
Q:=MM^{\mathsf T}.
\]
The established identity
\[
Q=A^3+2A^2+2I
\]
places the sector-overlap algebra inside the adjacency algebra of
\(L(\mathrm{Petersen})\).

The spectrum of \(Q\) is
\[
98^{(1)},\qquad
18^{(5)},\qquad
3^{(4)},\qquad
2^{(5)}.
\]
The eigenvalue \(98\) belongs to the one-dimensional all-ones
common-mode direction. Every centered direction has eigenvalue at
most \(18\).

After the canonical normalization
\[
\widehat Q:=\frac{1}{98}Q,
\]
the common mode is fixed while every centered relational direction
contracts. The sharp centered contraction factor is
\[
\frac{18}{98}=\frac{9}{49}.
\]
Thus
\[
\widehat Q^{\,k}\longrightarrow P_0,
\]
where \(P_0\) is the rank-one orthogonal projection onto the common
mode.

Yet every finite iterate
\[
\widehat Q^{\,k}
\]
is invertible.

This gives an exact finite linear realization of the distinction
between complete state and registered contrast:

\begin{quote}
A strong common aggregate can emerge while relative individuation
becomes arbitrarily weak, without any finite-stage merger of complete
linear states.
\end{quote}

We call this exact algebraic phenomenon
\emph{contrast concentration}.

The broader phrase
\emph{relational saturation} is reserved for a stronger,
projection-relative conjecture. Saturation requires not only a
contrast-concentrating operator, but also an admitted native loading
law, a declared registration, finite independent distinction capacity,
and retained complete information after the declared readout has
ceased to resolve additional independent contrast.

This paper proves contrast concentration.

It formulates relational saturation.

It does not prove that the native higher Thalean mechanics repeatedly
applies \(Q\), nor that physical systems saturate in this manner.
No identification is made here with physical gravity, black holes,
spacetime horizons, Planck-scale information density, or the cosmic
microwave background.

The theorem-side question therefore comes first:

\begin{quote}
Can bounded finite linear aggregation produce strong aggregate and
weak individuation at the same time?
\end{quote}

For the Thalean quadratic, the answer is exact.

\section{Scope and Imported Finite Structure}

\subsection{Purpose of the paper}

This paper gives a finite theorem of Thalean relational saturation on
the registered mechanics associated with
\[
\Gfifteen\cong L(\Petersen).
\]

The result combines three exact finite structures:

\begin{enumerate}[label=(\roman*)]

\item a native carrier language whose lawful completion family
contracts as
\[
8\longrightarrow2\longrightarrow1;
\]

\item a retained five-address-hidden sector selected exactly by the
context-activated native designation response;

\item the normalized Thalean quadratic, which independently contracts
centered registered contrast toward common mode.

\end{enumerate}

The paper does not introduce a new graph or alter the established
native carrier law.

It also does not identify native loading with iteration of the
Thalean quadratic.

That identification was an earlier conjectural possibility.  The
exact finite result is different and stronger: native selection and
quadratic contrast concentration are distinct certified operations
joined through a common finite response algebra.

\subsection{Imported theorem object}

We import the following finite results from the established
AT4val[60,6] theorem chain
\citep{cave2026dodecahedral,cave2026quantumrelativity}.

\begin{enumerate}[label=(\roman*)]

\item The quotient core is
\[
\Gfifteen\cong L(\Petersen).
\]

\item The transport-induced sector--edge incidence matrix
\[
M\in\{0,1\}^{15\times30}
\]
has row weight \(14\) and column weight \(7\).

\item Its Gram matrix is
\[
Q:=MM^{\mathsf T}
  =A^{3}+2A^{2}+2I,
\]
where \(A\) is the adjacency matrix of \(\Gfifteen\).

\item The adjacency spectrum is
\[
\spec(A)
=
\left\{
4^{(1)},\,
2^{(5)},\,
(-1)^{(4)},\,
(-2)^{(5)}
\right\},
\]
and therefore
\[
\spec(Q)
=
\left\{
98^{(1)},\,
18^{(5)},\,
3^{(4)},\,
2^{(5)}
\right\}.
\]

\item The signed lift
\[
\Gthirty\longrightarrow\Gfifteen
\]
contains nontrivial holonomy information not determined by \(Q\).

\item Declared projections in the higher finite construction can
return before the corresponding lifted or complete registered state
returns.

\item Registered history denotes retained finite information that
distinguishes complete states already identified by a declared
projection.

\item A declared real-valued register
\(h\in\mathbb R^n\) has centered relational magnitude
\[
R_n(h)
=
\left\|
h-\overline h\,\one_n
\right\|_2.
\]
This scalar is a many-to-one compression, not a complete state
coordinate.

\item The native registered carrier has \(60\) directed-edge
contextual states.  Every directed quotient edge has exactly two
registered continuations, one of designation \(b\) and one of
designation \(ab\), and every live oriented two-edge path has a unique
registered completion.

\item The accepted registered language contains \(120\) oriented
words.  Its exact completion profile is
\[
8\longrightarrow2\longrightarrow1\longrightarrow1\longrightarrow1.
\]

\item The five-address registration partitions the fifteen-state
sector register into a five-dimensional fibre-constant sector and a
ten-dimensional fibre-internal sector.

\end{enumerate}

These structures remain distinct.  In particular,
\[
\text{native selection}
\neq
\text{quadratic aggregation}
\neq
\text{covering holonomy}
\neq
\text{registered history}
\neq
\text{centered scalar magnitude}.
\]

\subsection{New contribution}

The upgraded theorem package has four parts.

First, for a live registered prefix \(p\), let
\[
\mathcal C(p)
\]
denote the lawful completion family.  Prefix extension is monotone:
\[
p\preceq p'
\quad\Longrightarrow\quad
\mathcal C(p')\subseteq\mathcal C(p).
\]
The certified native loading ladder is
\[
8\longrightarrow2\longrightarrow1.
\]
We therefore define native relational loading as accumulated
admissible constraint and registered-future saturation by
\[
|\mathcal C(p)|=1.
\]

Second, future uniqueness does not erase complete context.  The
two-edge saturated stage contains \(120\) distinct contexts with
\(120\) distinct registered completions, and the terminal completion
register retains two native holonomy designations over every closing
chamber.

Third, let \(\Delta\) denote the context-activated native
\(b/ab\) designation response on the fifteen-state chamber register.
Let \(H\) denote the exact ten-dimensional five-address-hidden
projector.  Then
\[
\Delta^{\mathsf T}\Delta
=
\Delta\Delta^{\mathsf T}
=
\frac34H.
\]
The range and domain of \(\Delta\) lie in the hidden sector, while
the five-address readout satisfies
\[
\rho\Delta=0.
\]

Fourth, the Thalean quadratic provides a separate
contrast-concentrating operation.  For
\[
\widehat Q:=\frac1{98}Q,
\]
the common mode is fixed and the complete centered subspace contracts
by the sharp operator factor \(9/49\).  Consequently,
\[
\widehat Q^{\,k}x
\longrightarrow
\Pzero x.
\]
Every finite iterate remains invertible.

The native loading response and the centered Thalean quadratic lie in
the same exact three-dimensional transport-response algebra, but the
two operators are not proportional and do not perform the same
operation.

\subsection{Native loading is not quadratic contrast concentration}

The distinction can be tested directly using the Paper 14 contrast
observable.

For the canonical uniform completion-occupancy register, native
loading gives
\[
\zeta:
\frac{31}{55}
\longrightarrow
\frac{41}{65}
\longrightarrow
\frac{7}{10}.
\]
The common component remains fixed while contrast strictly increases
on every one of the \(180\) strict loading extensions.

Thus
\[
\text{native relational loading}
\neq
\text{\(\widehat Q\)-type contrast concentration}.
\]

Native loading sharpens the selected future.

The normalized quadratic suppresses centered contrast relative to
common mode.

The theorem keeps these operations separate.

\subsection{Theorem boundary}

The following implications are not asserted:
\[
Q
=
\text{native temporal evolution},
\]
\[
\widehat Q
=
\text{native loading},
\]
\[
\widehat Q
=
\text{physical averaging law},
\]
or
\[
\text{relational saturation}
=
\text{physical gravitational collapse}.
\]

The exact theorem is finite and internal.

The native carrier, the retained-sector response, and the quadratic
contrast law are joined mathematically without identifying any of
them with a physical process.

\subsection{Three epistemic layers}

The paper maintains three layers.

\begin{description}

\item[Established finite mathematics.]
The graph, quotient hierarchy, incidence matrix, quadratic identity,
spectra, native carrier language, completion-family loading order,
context activation, projection-relative saturation, retained hidden
sector, covering receipts, recurrence distinctions, registered
history, Thalean time, and centered magnitude.

\item[Internal conditional mechanics.]
Clean release, overcompression, any proposed composition of native
selection with contrast concentration into a larger cycle, and any
stronger global inter-face transport construction.

\item[Physical bridge.]
Any proposed relation to physical information density, gravity,
black holes, horizons, singularities, the CMB, or Planck-scale
support.

\end{description}

No result is promoted from one layer to the next without an explicit
crosswalk.

\section{Projection and Retained Distinction}

\subsection{Declared projections}

\begin{definition}[Complete and projected state]
Let
\[
q:X\longrightarrow Y
\]
be a declared map between finite state spaces.  Relative to this map,
\(x\in X\) is the complete state and \(q(x)\in Y\) is its projected
state.
\end{definition}

The words complete and projected are map-relative.  They do not
assert that \(X\) is an absolute final ontology or that \(Y\) is the
only available registration.

\begin{definition}[Projection equivalence]
For \(x,x'\in X\), define
\[
x\sim_q x'
\quad\Longleftrightarrow\quad
q(x)=q(x').
\]
The equivalence classes are the fibres of \(q\).
\end{definition}

\begin{definition}[Projection-relative degeneracy]
A pair \(x\neq x'\) is projection-degenerate relative to \(q\) when
\[
q(x)=q(x').
\]
The complete distinction remains present in \(X\), while the declared
projection fails to resolve it.
\end{definition}

This is the exact finite meaning used in the present paper.  It is
not an assertion that two complete states have merged.

\subsection{Projected and complete return}

Let
\[
T:X\longrightarrow X
\]
be a declared finite transformation.

\begin{definition}[Return predicates]
Complete return at step \(m\) means
\[
T^{m}x=x.
\]
Projected return at step \(m\) means
\[
q(T^{m}x)=q(x).
\]
\end{definition}

\begin{proposition}[Registered-return implication]
For every declared finite transformation and projection,
\[
T^{m}x=x
\quad\Longrightarrow\quad
q(T^{m}x)=q(x).
\]
The converse need not hold.
\end{proposition}

\begin{proof}
The implication follows by applying \(q\) to the equality
\(T^{m}x=x\).  The failure of the converse is realized explicitly by
the certified lifted and registered recurrence examples imported from
the earlier theorem chain.
\end{proof}

Two established examples are:

\begin{enumerate}[label=(\roman*)]
\item a three-step projected triangle return whose twisted lift
requires six steps for complete return;

\item a ten-step quotient return in the registered five-frame
construction whose complete registered return requires twenty steps.
\end{enumerate}

These examples act on different finite objects and have different
receipt groups.  Only the structural implication is shared.

\subsection{Registration data and registered history}

\begin{definition}[Registration datum]
For a declared projection \(q:X\to Y\), a registration datum is
finite information retained in the complete construction that is not
determined by the projected state \(q(x)\).
\end{definition}

\begin{definition}[Registered history]
A registered history is an ordered finite composition of declared
transformations together with the retained registration data required
to distinguish the resulting complete state from its projected state.
\end{definition}

The established construction contains several examples:

\begin{itemize}
\item deck receipts in finite covers;
\item binary closure data;
\item binary edge gauge classes;
\item integer refinements;
\item registered frame coordinates;
\item primitive registered-history cycle classes.
\end{itemize}

These are not one universal receipt group.  Their common role is that
they preserve distinctions erased by a specified forgetful map.

\subsection{Scalar compression is another projection}

Let
\[
\rho:S\longrightarrow\mathbb R^n
\]
be a declared real-valued readout from a richer state space \(S\).
The composition
\[
S
\overset{\rho}{\longrightarrow}
\mathbb R^n
\overset{R_n}{\longrightarrow}
\mathbb R_{\ge0}
\]
is generally many-to-one twice over.

Distinct complete states may have the same vector readout.  Distinct
centered vectors may also have the same centered magnitude.  Thus
\[
R_n(\rho(s))=R_n(\rho(s'))
\]
does not imply
\[
s=s'.
\]

This provides an exact finite model for loss of registered
individuation without loss of complete distinction.

\subsection{No universal information ladder}

The finite constructions do not supply one canonical chain
\[
\text{complete state}
\longrightarrow
\text{lift}
\longrightarrow
\text{registration}
\longrightarrow
\text{scalar}
\]
covering every object in the theory.  They supply several declared
forgetful maps with separately certified domains.

Accordingly, every saturation claim in this paper must state:

\begin{enumerate}[label=(\roman*)]
\item the complete state space;
\item the declared projection or readout;
\item the distinction being retained;
\item the distinction being forgotten.
\end{enumerate}

\section{The Thalean Quadratic}

\subsection{The Petersen-rooted sector system}

Let
\[
\Gfifteen\cong L(\Petersen)
\]
and let \(A\) be its adjacency matrix.  The graph has fifteen
vertices and thirty edges.

For each \(v\in V(\Gfifteen)\), the transport construction supplies a
sector
\[
S(v)\subseteq E(\Gfifteen).
\]
Every sector contains fourteen edges, and every core edge belongs to
exactly seven sectors.

Define
\[
M=(M_{v,e})\in\{0,1\}^{15\times30}
\]
by
\[
M_{v,e}
=
\begin{cases}
1,&e\in S(v),\\
0,&e\notin S(v).
\end{cases}
\]

The Thalean quadratic is
\[
Q:=MM^{\mathsf T}.
\]
Its entries are
\[
Q_{u,v}
=
|S(u)\cap S(v)|.
\]

\subsection{Distance and adjacency forms}

The overlap depends only on distance in \(\Gfifteen\):
\[
Q
=
14D_0+9D_1+5D_2+4D_3.
\]
Because \(\Gfifteen\) is distance-regular of diameter three, this
reduces exactly to
\[
Q=A^3+2A^2+2I.
\]

The identity has two logically separate interpretations.

First, it is a Gram identity:
\[
Q=MM^{\mathsf T}.
\]
Second, it belongs to the adjacency algebra:
\[
Q=A^3+2A^2+2I.
\]

Neither description includes the signed lift or its holonomy class.

\subsection{Edge-mediated accumulation}

For a sector coefficient vector
\[
x\in\mathbb R^{15},
\]
the vector
\[
M^{\mathsf T}x\in\mathbb R^{30}
\]
assigns to each core edge the sum of the coefficients of all sectors
containing that edge.

Applying \(M\) again returns the accumulated edge values to the
sector register:
\[
x
\longmapsto
M^{\mathsf T}x
\longmapsto
MM^{\mathsf T}x
=
Qx.
\]

Thus \(Q\) is an exact finite local-to-global-to-local incidence sum.

\subsection{The common incidence mode}

Let \(\one_{15}\) and \(\one_{30}\) denote the all-ones vectors in
their respective dimensions.  Column weight seven gives
\[
M^{\mathsf T}\one_{15}
=
7\one_{30}.
\]
Row weight fourteen gives
\[
M\one_{30}
=
14\one_{15}.
\]
Consequently,
\[
Q\one_{15}
=
M(7\one_{30})
=
98\one_{15}.
\]

The number \(98=7\cdot14\) is therefore obtained directly from the
balanced incidence structure.

\subsection{Spectrum}

The adjacency spectrum is
\[
\spec(A)
=
\left\{
4^{(1)},\,
2^{(5)},\,
(-1)^{(4)},\,
(-2)^{(5)}
\right\}.
\]
For
\[
p(t)=t^3+2t^2+2,
\]
we have
\[
p(4)=98,\qquad
p(2)=18,\qquad
p(-1)=3,\qquad
p(-2)=2.
\]
Hence
\[
\spec(Q)
=
\left\{
98^{(1)},\,
18^{(5)},\,
3^{(4)},\,
2^{(5)}
\right\}.
\]

In particular, \(Q\) is symmetric positive definite.

\subsection{Canonical normalization}

Define
\[
\widehat{Q}
:=
\frac{1}{98}Q.
\]
Then
\[
\widehat{Q}\one_{15}
=
\one_{15}.
\]

Since \(Q\) is symmetric and every row sum is \(98\),
\(\widehat{Q}\) is symmetric and doubly stochastic.  No probabilistic
interpretation is required.  Here it is used only as the unique scalar
normalization that fixes the common incidence mode.

Its spectrum is
\[
\spec(\widehat{Q})
=
\left\{
1^{(1)},\,
\left(\frac{9}{49}\right)^{(5)},\,
\left(\frac{3}{98}\right)^{(4)},\,
\left(\frac{1}{49}\right)^{(5)}
\right\}.
\]

\subsection{Operator boundary}

The map \(\widehat{Q}\) is a canonically normalized incidence
operator.  The present theorem does not establish that one
application of \(\widehat{Q}\) is:

\begin{itemize}
\item one native transport step;
\item one Thalean-time increment;
\item one physical averaging event;
\item or one stage of physical compression.
\end{itemize}

Iteration is studied algebraically.  Its promotion to native loading
requires a separate crosswalk.

\section{Common Mode and Centered Relational Contrast}

\subsection{Orthogonal projections}

Let
\[
J=\one_{15}\one_{15}^{\mathsf T}.
\]
Define
\[
\Pzero
=
\frac{1}{15}J
\]
and
\[
\Pperp
=
I-\Pzero.
\]

Then \(\Pzero\) is the orthogonal projection onto
\[
\operatorname{span}\{\one_{15}\},
\]
and \(\Pperp\) is the orthogonal projection onto
\[
\one_{15}^{\perp}
=
\left\{
x\in\mathbb R^{15}:
\one_{15}^{\mathsf T}x=0
\right\}.
\]

Every \(x\in\mathbb R^{15}\) has the unique decomposition
\[
x
=
\Pzero x+\Pperp x.
\]

Writing
\[
\overline{x}
=
\frac{1}{15}\one_{15}^{\mathsf T}x,
\]
we have
\[
\Pzero x
=
\overline{x}\,\one_{15}.
\]

\subsection{Common and relational magnitudes}

Define
\[
\Ccommon(x)
:=
\|\Pzero x\|_2
=
\sqrt{15}\,|\overline{x}|
\]
and
\[
\Rrel(x)
:=
\|\Pperp x\|_2
=
\left\|
x-\overline{x}\,\one_{15}
\right\|_2.
\]

The second quantity is exactly the centered relational magnitude
specialized to the fifteen-sector register.

Orthogonality gives
\[
\|x\|_2^2
=
\Ccommon(x)^2+\Rrel(x)^2.
\]

\subsection{Dimensionless contrast fraction}

For \(x\neq0\), define
\[
\zetaRel(x)
:=
\frac{\Rrel(x)^2}{\|x\|_2^2}.
\]
Then
\[
0\le\zetaRel(x)\le1.
\]

Moreover,
\[
\zetaRel(x)=0
\quad\Longleftrightarrow\quad
x\in\operatorname{span}\{\one_{15}\},
\]
while
\[
\zetaRel(x)=1
\quad\Longleftrightarrow\quad
x\in\one_{15}^{\perp}.
\]

Thus \(\zetaRel\) measures the fraction of Euclidean register
magnitude carried by centered relational contrast rather than by the
common mode.

When \(\Ccommon(x)\neq0\), we also define the relative contrast ratio
\[
\rho(x)
:=
\frac{\Rrel(x)}{\Ccommon(x)}.
\]

\subsection{Invariant decomposition under the quadratic}

Because \(Q\) is a polynomial in the regular graph adjacency matrix
\(A\), both \(\Pzero\) and \(\Pperp\) commute with \(Q\), and hence
with \(\widehat{Q}\):
\[
\Pzero\widehat{Q}
=
\widehat{Q}\Pzero,
\qquad
\Pperp\widehat{Q}
=
\widehat{Q}\Pperp.
\]

Therefore the common and centered subspaces are invariant.

This point is essential.  Contrast concentration under
\(\widehat{Q}\) does not arise because centered amplitude is
transferred into the common mode.  The two subspaces do not mix.
Instead, the common mode has eigenvalue \(1\), while every centered
mode has eigenvalue strictly smaller than \(1\).

\subsection{Complete spectral decomposition}

Let \(E_\lambda\) denote the orthogonal eigenspace of \(A\) with
eigenvalue \(\lambda\).  Then
\[
\mathbb R^{15}
=
E_4
\oplus
E_2
\oplus
E_{-1}
\oplus
E_{-2},
\]
with
\[
\dim E_4=1,\quad
\dim E_2=5,\quad
\dim E_{-1}=4,\quad
\dim E_{-2}=5.
\]

The common-mode space is
\[
E_4=\operatorname{span}\{\one_{15}\}.
\]

On these four spaces, \(\widehat{Q}\) acts respectively by
\[
1,\qquad
\frac{9}{49},\qquad
\frac{3}{98},\qquad
\frac{1}{49}.
\]

The largest centered eigenvalue is therefore
\[
\lambda_{\mathrm{cent}}
=
\frac{9}{49}.
\]

\section{The Contrast-Concentration Theorem}

\begin{theorem}[Thalean quadratic contrast-concentration theorem]
Let
\[
\widehat{Q}
=
\frac{1}{98}
\left(A^3+2A^2+2I\right)
\]
on
\[
\mathbb R^{15},
\]
where \(A\) is the adjacency matrix of
\(\Gfifteen\cong L(\Petersen)\).

For every \(x\in\mathbb R^{15}\) and every integer \(k\ge0\):

\begin{enumerate}[label=(\roman*)]
\item the common component is fixed,
\[
\Pzero\widehat{Q}^{\,k}x
=
\Pzero x;
\]

\item the centered component satisfies
\[
\Rrel(\widehat{Q}^{\,k}x)
\le
\left(\frac{9}{49}\right)^k
\Rrel(x);
\]

\item the iterates converge to the common-mode projection,
\[
\lim_{k\to\infty}
\widehat{Q}^{\,k}x
=
\Pzero x
=
\overline{x}\,\one_{15};
\]

\item if \(\Ccommon(x)\neq0\), then
\[
\rho(\widehat{Q}^{\,k}x)
\le
\left(\frac{9}{49}\right)^k
\rho(x);
\]

\item if \(x\neq0\) and \(\Ccommon(x)\neq0\), then
\[
\zetaRel(\widehat{Q}^{\,k}x)
\le
\left(\frac{9}{49}\right)^{2k}
\rho(x)^2.
\]
\end{enumerate}
\end{theorem}

\begin{proof}
The matrix \(A\) is real symmetric, so its eigenspaces form an
orthogonal decomposition.  Since
\[
Q=A^3+2A^2+2I,
\]
the matrices \(A\), \(Q\), and \(\widehat{Q}\) have the same
eigenspaces.

On the all-ones eigenspace \(E_4\), the operator
\(\widehat{Q}\) has eigenvalue \(1\).  Therefore
\[
\Pzero\widehat{Q}^{\,k}x
=
\Pzero x.
\]

On the centered subspace
\[
\one_{15}^{\perp}
=
E_2\oplus E_{-1}\oplus E_{-2},
\]
the eigenvalues are
\[
\frac{9}{49},
\qquad
\frac{3}{98},
\qquad
\frac{1}{49}.
\]
All are nonnegative, and their maximum is \(9/49\).  Hence the
Euclidean operator norm of
\(\widehat{Q}^{\,k}\) restricted to the centered subspace is
\[
\left(\frac{9}{49}\right)^k.
\]
Thus
\[
\begin{aligned}
\Rrel(\widehat{Q}^{\,k}x)
&=
\left\|
\Pperp\widehat{Q}^{\,k}x
\right\|_2\\
&=
\left\|
\widehat{Q}^{\,k}\Pperp x
\right\|_2\\
&\le
\left(\frac{9}{49}\right)^k
\|\Pperp x\|_2.
\end{aligned}
\]

The centered bound tends to zero, while the common component remains
fixed.  Therefore
\[
\widehat{Q}^{\,k}x
=
\Pzero x+\Pperp\widehat{Q}^{\,k}x
\longrightarrow
\Pzero x.
\]

When \(\Ccommon(x)\neq0\), common-mode invariance gives
\[
\Ccommon(\widehat{Q}^{\,k}x)
=
\Ccommon(x),
\]
so division yields the bound for \(\rho\).

Finally,
\[
\zetaRel(y)
=
\frac{\Rrel(y)^2}
{\Ccommon(y)^2+\Rrel(y)^2}
\le
\frac{\Rrel(y)^2}{\Ccommon(y)^2}.
\]
Applying the preceding two bounds with
\(y=\widehat{Q}^{\,k}x\) proves the final inequality.
\end{proof}

\subsection{Exact modal formula}

Write
\[
x
=
x_4+x_2+x_{-1}+x_{-2},
\qquad
x_\lambda\in E_\lambda.
\]
Then
\[
\widehat{Q}^{\,k}x
=
x_4
+
\left(\frac{9}{49}\right)^k x_2
+
\left(\frac{3}{98}\right)^k x_{-1}
+
\left(\frac{1}{49}\right)^k x_{-2}.
\]

Consequently,
\[
\Rrel(\widehat{Q}^{\,k}x)^2
=
\left(\frac{9}{49}\right)^{2k}\|x_2\|_2^2
+
\left(\frac{3}{98}\right)^{2k}\|x_{-1}\|_2^2
+
\left(\frac{1}{49}\right)^{2k}\|x_{-2}\|_2^2.
\]

The bound in the theorem is sharp whenever the centered component
lies entirely in \(E_2\).

\begin{corollary}[Operator convergence]
In Euclidean operator norm,
\[
\widehat{Q}^{\,k}
\longrightarrow
\Pzero,
\]
and
\[
\left\|
\widehat{Q}^{\,k}-\Pzero
\right\|_{2\to2}
=
\left(\frac{9}{49}\right)^k.
\]
\end{corollary}

\begin{proof}
The difference vanishes on \(E_4\) and acts on the centered
eigenspaces with eigenvalues
\[
\left(\frac{9}{49}\right)^k,\quad
\left(\frac{3}{98}\right)^k,\quad
\left(\frac{1}{49}\right)^k.
\]
The largest absolute value is
\((9/49)^k\).
\end{proof}

\begin{remark}[Exactly centered input]
If
\[
\Pzero x=0,
\]
then every iterate remains centered:
\[
\Pzero\widehat{Q}^{\,k}x=0.
\]
In that case,
\[
\widehat{Q}^{\,k}x\longrightarrow0.
\]
The operator does not create a common component that was absent from
the input.
\end{remark}

\begin{remark}[Raw and normalized aggregation]
For the unnormalized quadratic,
\[
\Ccommon(Q^kx)
=
98^k\Ccommon(x)
\]
and
\[
\Rrel(Q^kx)
\le
18^k\Rrel(x).
\]
Thus the same ratio
\[
\frac{18}{98}
=
\frac{9}{49}
\]
governs relative contrast concentration under raw \(Q\)-iteration.

The normalized operator is preferable for the theorem because it
fixes the aggregate scale and turns the relative statement into an
absolute centered contraction.
\end{remark}

\subsection{Interpretive boundary}


The theorem proves exact algebraic convergence under repeated
application of a canonically normalized incidence operator.

It does not assert that native Thalean loading repeatedly applies
\(\widehat Q\), and it does not identify \(\widehat Q\) with physical
dynamics.

The native theorem now supplies a different exact operation:
\[
8\longrightarrow2\longrightarrow1
\]
for lawful completion-family contraction.

Under the canonical completion-occupancy register, native loading
strictly increases Paper 14 contrast on every strict loading
extension.

Therefore the former candidate identification
\[
\text{native loading}
\longrightarrow
\widehat Q\text{-type concentration}
\]
is rejected.

The correct relationship is
\[
\text{native selection}
\neq
\text{quadratic concentration},
\]
with both operations joined through the certified finite
transport-response algebra.

\section{Invertibility and Non-Erasure}

\subsection{Finite-stage invertibility}

All eigenvalues of \(Q\) are strictly positive:
\[
98,\quad18,\quad3,\quad2.
\]
Therefore \(Q\) and \(\widehat{Q}\) are invertible.

Their determinants are
\[
\det Q
=
98\cdot18^5\cdot3^4\cdot2^5
>
0
\]
and
\[
\det\widehat{Q}
=
\left(\frac{9}{49}\right)^5
\left(\frac{3}{98}\right)^4
\left(\frac{1}{49}\right)^5
>
0.
\]

\begin{proposition}[Finite-stage non-erasure]
For every finite integer \(k\ge0\) and every
\(x,y\in\mathbb R^{15}\),
\[
\widehat{Q}^{\,k}x
=
\widehat{Q}^{\,k}y
\quad\Longleftrightarrow\quad
x=y.
\]
\end{proposition}

\begin{proof}
Every finite power of an invertible matrix is invertible.
\end{proof}

Thus no two complete linear sector states are merged by a finite
number of normalized quadratic aggregation steps.

\subsection{An invertible sequence with a noninvertible limit}

The contrast-concentration theorem gives
\[
\widehat{Q}^{\,k}
\longrightarrow
\Pzero.
\]
The limiting operator \(\Pzero\) has rank one and is not invertible.

Therefore Paper 14 exhibits the following exact finite-dimensional
phenomenon:
\[
\begin{array}{c}
\text{every finite iterate is invertible},\\[0.3em]
\text{the operator limit is many-to-one}.
\end{array}
\]

There is no contradiction.  The inverses
\[
\widehat{Q}^{-k}
\]
become increasingly ill-conditioned on centered directions as
\(k\) grows.

The finite state remains mathematically recoverable when the exact
iterate and exact output are retained.  Recovery becomes
progressively sensitive to any loss of centered precision.

\subsection{Exact distinction and effective resolution}

Let
\[
d=x-y\neq0
\]
be a complete distinction vector.  Then
\[
\widehat{Q}^{\,k}d\neq0
\]
for every finite \(k\).

If \(d\) is centered, however,
\[
\left\|
\widehat{Q}^{\,k}d
\right\|_2
\le
\left(\frac{9}{49}\right)^k
\|d\|_2.
\]

Thus a complete distinction may remain exact while its centered
registered amplitude becomes arbitrarily small.

For a declared resolution threshold
\(\varepsilon>0\), there may exist a finite \(k\) for which
\[
0
<
\left\|
\widehat{Q}^{\,k}d
\right\|_2
<
\varepsilon.
\]
The distinction is still present in exact arithmetic but is no
longer resolved by that thresholded readout.

This motivates the phrase
\emph{effective registration degeneracy}.  It must not be confused
with exact equality of complete states.

\subsection{Scalar readout loss}

The centered scalar
\[
R_{15}(x)
=
\|\Pperp x\|_2
\]
is itself many-to-one.  Distinct centered vectors may share the same
norm.  Therefore even before any iteration,
\[
R_{15}(x)=R_{15}(y)
\]
does not imply
\[
x=y.
\]

The normalized contrast fraction
\[
\zetaRel(x)
=
\frac{R_{15}(x)^2}{\|x\|_2^2}
\]
forgets still more structure: common scale, centered direction, and
all lifted data absent from the fifteen-component register.

Consequently, three levels must remain separate:

\begin{enumerate}[label=(\roman*)]
\item exact complete sector state;
\item fifteen-component incidence register;
\item scalar contrast readout.
\end{enumerate}

Loss at level (iii) is not erasure at level (i).

\subsection{Relation to retained history}

The earlier finite mechanics supplies stronger distinctions than the
quadratic sector register alone.  Covering holonomy, deck receipts,
integer refinements, and registered-history winding can distinguish
states or histories that have already become identical under a
coarser projection.

The present quadratic theorem does not yet act on those richer
objects.  It supplies the contrast-concentrating side of the proposed
saturation architecture.  The retained-history theorem supplies
examples of complete distinction surviving projection.

A full internal saturation theorem would require one admitted native
construction in which both occur in the same declared diagram:
\[
\begin{array}{ccc}
\text{complete admitted state}
&\longrightarrow&
\text{retained receipt or winding}\\
\downarrow
&&
\downarrow\\
\text{contrast-concentrating register}
&\longrightarrow&
\text{coarse scalar or visible state}.
\end{array}
\]

The current theorem chain has not yet supplied that crosswalk.

\subsection{The exact conclusion}

The exact conclusion of Sections 4--7 is:

\begin{quote}
The normalized Thalean quadratic contracts every centered sector
direction by a factor at most \(9/49\) per iterate, converges to the
common-mode projection, and remains invertible at every finite
iterate.
\end{quote}

This is a theorem of finite linear algebra.

Its saturation interpretation is internal and conditional.

Its physical interpretation remains open.

\section{Projection-Relative Saturation}

\subsection{Why saturation must be projection-relative}

The finite mechanics contains several distinct forgetful maps.
Consequently, the statement that a system has ``lost information''
is incomplete unless the relevant map is named.

Let
\[
q:X\longrightarrow Y
\]
be a declared projection.

A distinction between \(x,x'\in X\) survives the projection when
\[
q(x)\neq q(x').
\]
It is hidden by the projection when
\[
x\neq x',
\qquad
q(x)=q(x').
\]

The second relation does not imply that the complete states have
merged. It says only that the declared projection no longer resolves
their difference.

\begin{definition}[Projection-relative degeneracy]
Two distinct complete states \(x,x'\in X\) are projection-degenerate
relative to \(q\) when
\[
q(x)=q(x').
\]
\end{definition}

This is an exact statement about a map.

It is not an ontological identification of \(x\) and \(x'\).

\subsection{Independent registered distinction}

Let
\[
\mathcal X\subseteq X
\]
be a declared family of complete states.

The projected state set is
\[
q(\mathcal X)\subseteq Y.
\]

When the relevant notion of independence is cardinal, the number of
projected distinction classes is simply
\[
K_q(\mathcal X)
=
|q(\mathcal X)|.
\]

More generally, when a declared linear or combinatorial rank function
is available, write
\[
r_q(\mathcal X)
\]
for the independent distinction rank that survives the projection.

No universal rank function is assumed in this paper.

The choice must be stated for the particular register being studied.

\subsection{Projection-relative saturation}

\begin{definition}[Projection-relative saturation]
Let
\[
\{x_\lambda\}_{\lambda\in\Lambda}
\]
be a declared loading family of complete states or admitted histories,
and let
\[
q:X\longrightarrow Y
\]
be a declared projection or registration.

The family enters a projection-relative saturation regime when:

\begin{enumerate}[label=(\roman*)]
\item complete distinction remains nontrivial;

\item the independent distinction visible through \(q\) is bounded;

\item increasing loading ceases to produce corresponding growth of
independently resolved projected distinction;

\item additional complete variation is therefore increasingly carried
inside fibres of \(q\), retained receipts, lifted structure, or other
higher-resolution data.
\end{enumerate}
\end{definition}

Saturation is therefore not defined as disappearance of the complete
state.

It is a statement about the relationship between complete variation
and the resolving capacity of a declared registration.

\subsection{Contrast saturation on the sector register}

For the fifteen-sector real register, Paper 14 defines
\[
\zetaRel(x)
=
\frac{\|\Pperp x\|_2^2}{\|x\|_2^2},
\qquad
x\neq0.
\]

This dimensionless quantity measures the fraction of register
magnitude carried by centered relational contrast.

\begin{definition}[Contrast-saturating sequence]
A sequence of nonzero registered states
\[
x_0,x_1,x_2,\ldots
\]
is contrast-saturating when
\[
\zetaRel(x_k)
\]
approaches a stable minimum while the corresponding complete state
sequence remains nontrivial.
\end{definition}

For normalized quadratic aggregation, if
\[
\Pzero x\neq0,
\]
the contrast-concentration theorem gives
\[
\zetaRel(\widehat Q^{\,k}x)
\longrightarrow0.
\]

Thus \(\widehat Q\)-iteration supplies an exact algebraic example of
complete common-mode concentration.

\begin{corollary}[Asymptotic contrast saturation]
If
\[
\Pzero x\neq0,
\]
then
\[
\lim_{k\to\infty}
\zetaRel(\widehat Q^{\,k}x)
=
0.
\]
For every finite \(k\), however,
\[
\widehat Q^{\,k}
\]
remains invertible.
\end{corollary}

Hence the same finite sequence exhibits
\[
\text{relative registered contrast}\longrightarrow0
\]
while
\[
\text{exact finite-stage distinction remains recoverable}.
\]

This is the exact algebraic core motivating the saturation language.

\subsection{Finite-resolution saturation}

Exact arithmetic is not the only possible declared registration.

Let
\[
\varepsilon>0
\]
be a fixed resolution threshold, and let
\[
d=x-y
\]
be a nonzero centered distinction:
\[
d\in\one_{15}^{\perp}.
\]

The contrast-concentration theorem gives
\[
\|\widehat Q^{\,k}d\|_2
\le
\left(\frac{9}{49}\right)^k
\|d\|_2.
\]

Therefore, for every fixed \(\varepsilon>0\), there exists a finite
\(k_\varepsilon\) such that
\[
0<
\|\widehat Q^{\,k}d\|_2
<
\varepsilon
\]
for all
\[
k\ge k_\varepsilon.
\]

The distinction remains mathematically nonzero, but a registration
with resolution \(\varepsilon\) no longer distinguishes it.

We call this
\emph{effective registration degeneracy}.

It must remain distinct from exact equality.

\subsection{Saturation is not global incompatibility}

The conjecture presupposes an admitted globally compatible transport
structure.

Failure of locally admissible choices to assemble into a globally
closed transport relation is a different phenomenon.

We therefore distinguish:

\begin{description}

\item[Local inadmissibility.]
A proposed local transition is not contained in the admitted local
transport vocabulary.

\item[Global incompatibility.]
Locally admissible transitions cannot be assembled into a globally
closed transport assignment satisfying the declared compatibility
conditions.

\item[Relational saturation.]
A globally admitted construction retains complete distinctions while
a declared registration ceases gaining independent resolving power.

\end{description}

The third must never be inferred from the second.

\subsection{Boundary}

The exact result of this section is a definition plus a proved
algebraic example.

Normalized \(Q\)-iteration is contrast-saturating relative to the
centered sector readout.

What remains unproved is that a native admitted Thalean loading law
realizes this sequence or an equivalent contrast-concentrating family.

Accordingly,
\[
\text{contrast saturation under }\widehat Q
\]
is theorem-side, while
\[
\text{native relational saturation}
\]
remains conjectural.

\section{Three Finite Summation Laws}

The established finite mechanics contains several exact additive or
accumulative structures. Together they support a general architectural
principle:

\begin{quote}
Global relational quantities may arise through lawful sums of finite
local or compositional contributions.
\end{quote}

The examples below act on different mathematical objects. They are
not identified with one another.

\subsection{Incidence accumulation}

Let
\[
x\in\mathbb R^{15}
\]
be a sector coefficient vector.

The transport-induced sector-edge incidence matrix gives
\[
x
\longmapsto
M^{\mathsf T}x
\longmapsto
MM^{\mathsf T}x
=
Qx.
\]

The intermediate vector
\[
M^{\mathsf T}x\in\mathbb R^{30}
\]
assigns to each core edge the sum of the sector coefficients incident
through the transport-induced sector system.

Applying \(M\) returns those accumulated edge values to the
fifteen-sector register.

Thus
\[
Q=MM^{\mathsf T}
\]
is an exact finite local-to-global-to-local incidence sum.

Because
\[
Q
=
14D_0+9D_1+5D_2+4D_3,
\]
the component at a sector basepoint \(v\) is
\[
(Qx)_v
=
14x_v
+
9\sum_{d(v,u)=1}x_u
+
5\sum_{d(v,u)=2}x_u
+
4\sum_{d(v,u)=3}x_u.
\]

The coefficients
\[
14,\quad9,\quad5,\quad4
\]
are exact sector-overlap multiplicities.

They are not fitted long-range weights.

\subsection{Balanced incidence and common aggregation}

Every row of \(M\) has weight \(14\), and every column has weight
\(7\). Therefore
\[
M\one_{30}
=
14\one_{15}
\]
and
\[
M^{\mathsf T}\one_{15}
=
7\one_{30}.
\]

Consequently,
\[
Q\one_{15}
=
98\one_{15}.
\]

The common aggregate is therefore already determined by the balanced
finite incidence structure:
\[
98=14\cdot7.
\]

This is the common mode isolated by the contrast-concentration
theorem.

\subsection{Binary regional composition}

The twelve-section closure construction contains a binary edge
cochain
\[
\delta
\]
and triangle charge
\[
\omega
\]
satisfying
\[
d\delta=\omega.
\]

For any declared mod-two region \(R\) assembled from native
triangles, interior edges occur twice and cancel. Hence
\[
\bigoplus_{\tau\subset R}\omega(\tau)
=
\bigoplus_{e\subset\partial R}\delta(e).
\]

This is a finite cochain identity.

It expresses a global regional receipt as a sum of boundary
contributions.

No continuous differential geometry is required.

\subsection{Integer regional composition}

The same binary closure class admits certified integer parents.

Let
\[
\eta
\]
denote either admitted parent and define
\[
\Omega=d\eta.
\]

For an oriented finite triangle region \(R\),
\[
\sum_{\tau\subset R}\Omega(\tau)
=
\sum_{e\subset\partial R}\eta(e).
\]

Interior oriented edge contributions cancel.

Reduction modulo two recovers the binary law:
\[
\Omega\bmod2=\omega.
\]

Thus the integer regional sum contains signed additive information
that is absent from the binary reduction.

The binary and integer laws are related by an explicit reduction map,
but they are not the same mathematical object.

\subsection{Registered-history winding}

The certified registered-history construction supplies a third form
of finite accumulation.

Thalean time is
\[
\tauT(\gamma)
=
\operatorname{wind}_{r_1}(\gamma),
\]
where \(r_1\) is the unique \(G_{30}\)-visible primitive
registered-history cycle species.

For composable registered histories,
\[
\tauT(\gamma_2\circ\gamma_1)
=
\tauT(\gamma_2)
+
\tauT(\gamma_1).
\]

Under reversal,
\[
\tauT(\gamma^{-1})
=
-\tauT(\gamma).
\]

This is an additive history coordinate.

It is not an incidence sum on \(G_{15}\), and it is not the regional
integer charge on the twelve-section closure complex.

\subsection{Three domains of summation}

The three exact structures may therefore be summarized as

\[
x
\longmapsto
Qx
\]
on the sector register,

\[
\delta
\longmapsto
d\delta=\omega
\]
on the finite closure complex,

and

\[
\gamma
\longmapsto
\tauT(\gamma)
\]
on registered-history cycles.

Their common feature is lawful finite accumulation.

Their domains, codomains, coefficient systems, and provenance are
different.

No theorem in this paper collapses them into a universal sum.

\subsection{QR long range as accumulated local contribution}

The finite examples motivate the following internal architectural
principle:

\begin{quote}
QR long-range structure is sought as the accumulated result of
admissible finite contributions, not as an independently inserted
long-range selector.
\end{quote}

This principle does not imply that every locally admissible
contribution participates in the global sum.

The summation domain must itself be selected by the relevant
compatibility law.

Schematically,
\[
\text{local contribution}
\longrightarrow
\text{global compatibility}
\longrightarrow
\text{admitted summation domain}
\longrightarrow
\text{registered aggregate}.
\]

The unresolved inter-face transport problem discussed in the next
section lies precisely at the global-compatibility step.

\subsection{Boundary}

The exact theorem-side content of this section is limited to the
three established finite sums.

The broader claim that these structures are manifestations of one
physical long-range law is not made.

In particular,
\[
Qx,
\qquad
\sum\Omega,
\qquad
\tauT
\]
must remain distinct unless an explicit crosswalk is proved.

\section{Global Compatibility and the Inter-Face Sigma Frontier}

\subsection{Local law and global closure are distinct}

Project 41 also contains a stronger inter-face transport problem that
is logically separate from the native registered carrier used in the
relational-saturation theorem.

The seam construction supplies a locally admissible inter-face
transport vocabulary.

Within that construction, the local \(r_3\) action has a universal
form: it reverses side and polarity while preserving sheet.  Its
section routing contains both same-section and changed-section
branches.

The completed local vocabulary is also closed under the independently
derived neighbor-frame involution \(E\).

These are local and covariance results.

They do not by themselves construct a globally consistent seam
transport.

The remaining requirement for that stronger construction is that
selected seam actions satisfy the declared four-step closure relation
\[
(r_2r_3)^4=1
\]
throughout the complete flag system.

Thus the architecture continues to distinguish
\[
\text{local admissibility}
\]
from
\[
\text{global compatibility}.
\]

\subsection{The inter-face sigma}

\begin{definition}[Inter-face sigma]
An inter-face sigma is a globally selected assignment of locally
admissible seam transport relations across adjacent faces satisfying
all declared covariance and closure conditions.
\end{definition}

Sigma is not an additional long-range selector placed on top of the
local mechanics.

It is a compatibility relation for one particular stronger
inter-face transport construction.

Schematically,
\[
\text{local seam vocabulary}
\longrightarrow
\text{covariant compatibility}
\longrightarrow
\sigma
\longrightarrow
\text{admitted global inter-face transport}.
\]

\subsection{Completed E-closed seam vocabulary}

Neighbor-frame covariance enlarges the local admissible vocabulary.

After canonical completion under the involution \(E\), the global
search used
\[
2816
\]
distinct local seam options.

This completed vocabulary is the correct search domain for the
present global \(r_3\) problem.

The enlargement is mathematically important.

It shows that the earlier truncated vocabulary did not exhaust the
locally admissible seam possibilities once neighbor-frame covariance
was imposed.

But a larger locally admissible vocabulary does not imply global
closure.

\subsection{Completed Mac search}

The completed Mac solver run searched the \(2816\)-option
\(E\)-closed seam vocabulary for a globally consistent \(r_3\).

The run reached
\[
376
\]
total solver iterations and accumulated
\[
175069
\]
unique four-step closure nogoods.

No complete global \(r_3\) was found.

The search also did not prove the completed search space infeasible.

The result therefore remains logically open:
\[
\text{no solution found}
\]
and
\[
\text{no nonexistence theorem}.
\]

\subsection{Partial closure improvement}

The completed search nevertheless moved measurably deeper into the
closure constraints.

Earlier candidates typically failed the four-step relation on all
\[
1920
\]
flags.

Improved candidates first reached
\[
1904
\]
failed flags, corresponding to
\[
16
\]
correctly closing flags.

The completed Mac run reached repeated candidates with
\[
1888
\]
failed flags.

Thus the best observed partial closure increased to
\[
1920-1888=32
\]
correctly closing flags.

This does not constitute a global solution.

It does show that completion under neighbor-frame covariance admits
strictly deeper partial closure than the earlier truncated search.

\subsection{What the search establishes}

The Mac work supports the following internal conclusions:

\begin{enumerate}[label=(\roman*)]

\item local admissibility and global compatibility are distinct;

\item neighbor-frame covariance changes the admissible local search
space;

\item substantial local admissibility does not guarantee global
closure;

\item the unresolved global object is an inter-face transport
relation rather than an independent long-range selector;

\item the remaining obstruction concerns this stronger inter-face
construction, not the existence of the already certified registered
carrier.

\end{enumerate}

These conclusions constrain any future global transport extension.

They do not weaken the finite relational-saturation theorem proved on
the independently certified registered carrier.

\subsection{Sigma and a possible global summation domain}

Section 9 discusses lawful summation of finite contributions.

For a future inter-face construction, the sum cannot be taken over
arbitrary locally admissible seam terms.

Its admissible summation domain would have to be compatible with the
global inter-face transport relation.

Thus a possible future chain is
\[
\text{local seam contribution}
\longrightarrow
\sigma
\longrightarrow
\text{global inter-face summation domain}
\longrightarrow
\text{registered aggregate}.
\]

This remains a conditional extension.

It is not the mechanism used to prove native relational saturation in
Section 8.

\subsection{Circulation and winding}

If a globally compatible inter-face sigma is eventually constructed,
one may evaluate its accumulated transport around closed admissible
circuits.

A nontrivial closed-circuit accumulation would provide a natural
candidate for a discrete circulation or winding quantity.

That possibility remains conditional.

No theorem in this paper identifies unresolved sigma circulation with:

\begin{itemize}

\item the signed \(G_{30}\to G_{15}\) covering cocycle;

\item the twelve-section binary closure class;

\item either certified integer parent;

\item the primitive registered-history species \(r_1\);

\item or Thalean time
\[
\tau_{\mathrm T}=\operatorname{wind}_{r_1}.
\]

\end{itemize}

Those structures have different domains and independent provenance.

An explicit crosswalk would be required before any identification is
permitted.

\subsection{Compatibility failure is not saturation}

The distinction remains decisive.

A proposed global inter-face transport assignment may fail because no
collection of locally admissible seam choices satisfies all of its
closure constraints.

That is global incompatibility for that construction.

Relational saturation is different.  It concerns an already admitted
carrier or registration in which future freedom or projected
resolving power saturates while complete distinction remains.

Therefore
\[
\text{global incompatibility}
\neq
\text{relational saturation}.
\]

A failed closure search cannot be reinterpreted as evidence for
saturation.

Nor does a failed or unfinished sigma search invalidate saturation
proved on another admitted finite carrier.

\subsection{Relation to the finite saturation theorem}

Earlier versions of this paper treated a globally closed inter-face
sigma as a prerequisite for any native relational-saturation theorem.

That prerequisite is no longer required.

The finite theorem of Section 8 is carried by an independently
certified native registered carrier with exact successor,
completion, and retained-context laws.

Its saturation mechanism is
\[
8\longrightarrow2\longrightarrow1,
\]
together with retained distinction and
\[
\Delta^{\mathsf T}\Delta=\frac34H.
\]

The unresolved sigma problem therefore has a narrower status:

\begin{quote}
Sigma remains an open candidate for a stronger global inter-face
transport theory, not a missing premise of the native finite
relational-saturation theorem.
\end{quote}

\subsection{Current boundary}

At the present boundary:

\[
\text{native registered carrier}
\quad
\text{established},
\]

\[
\text{native finite relational saturation}
\quad
\text{established},
\]

\[
\text{local seam law}
\quad
\text{established},
\]

\[
E\text{-closed local vocabulary}
\quad
\text{established},
\]

\[
\text{deeper partial inter-face closure}
\quad
\text{observed},
\]

\[
\text{global inter-face sigma}
\quad
\text{open}.
\]

No physical consequence follows from the open sigma frontier.

\section{Clean Release and Overcompression}

The contrast-concentration theorem establishes an exact finite
mechanism by which common-mode dominance can increase relative to
centered relational contrast.

It does not by itself determine an optimal loading regime.

This section introduces two internal concepts motivated by that
theorem:

\[
\text{clean release}
\]
and
\[
\text{overcompression}.
\]

Neither concept is part of the theorem proved in Sections 4--7.

\subsection{A declared output register}

Let
\[
\rho:X\longrightarrow H
\]
be a declared finite real-valued output register, with
\[
H\subseteq\mathbb R^n.
\]

For
\[
h=\rho(x),
\]
write the orthogonal decomposition
\[
h
=
P_0h+P_\perp h.
\]

The common component is
\[
P_0h
=
\overline h\,\one_n,
\]
while the centered component is
\[
P_\perp h
=
h-\overline h\,\one_n.
\]

Define
\[
C(h)
=
\|P_0h\|_2
\]
and
\[
R(h)
=
\|P_\perp h\|_2.
\]

A strong registered output need not maximize either quantity
separately.

A pure common-mode output may have large total magnitude while
carrying no internal coordinate contrast.

Conversely, a highly contrasted state need not be organized into a
robust outward registration.

\subsection{Coherence and individuation are different}


For the present internal mechanics, we distinguish two descriptive
properties.

\paragraph{Coherent registration.}
Multiple admissible contributions reinforce a common registered output
rather than cancelling one another.

\paragraph{Registered individuation.}
The declared output retains enough centered distinction to resolve the
relational alternatives relevant to that registration.

These properties need not vary together.

The upgraded native theorem now places an important restriction on any
proposed clean-release cycle.

Native relational loading itself does not suppress Paper 14 contrast
under the canonical completion register.

Instead,
\[
8\longrightarrow2\longrightarrow1
\]
sharpens the selected future and gives
\[
\zeta:
\frac{31}{55}
\longrightarrow
\frac{41}{65}
\longrightarrow
\frac{7}{10}.
\]

Thus any later loss of centered contrast requires a distinct
contrast-concentrating operation, such as the normalized quadratic,
or an explicitly proved composition law joining selection and
concentration.

Native selection and contrast suppression must not be silently
identified.

\subsection{Clean release}

\begin{definition}[Clean release]
A clean-release regime is a loading regime in which increasing
relational organization improves the coherence, robustness, or
magnitude of a declared outward registration while preserving the
independent contrast required by that output.
\end{definition}

The definition is projection-relative.

It does not assert a universal scalar measure of release quality.

A quantitative clean-release criterion would require at least:

\begin{enumerate}[label=(\roman*)]

\item a declared loading parameter;

\item a declared outward registration;

\item a measure of coherent registered output;

\item a measure of surviving independent contrast.

\end{enumerate}

No unique combination of these quantities is selected in this paper.

\subsection{Why maximum aggregate need not be maximum useful output}

Suppose a loading family is indexed by
\[
\lambda.
\]

Let
\[
C(\lambda)
\]
measure coherent outward registration and let
\[
D(\lambda)
\]
measure the independent contrast retained by that registration.

A useful output statistic
\[
U(\lambda)
\]
may depend positively on both.

For example, schematically,
\[
U(\lambda)
=
F(C(\lambda),D(\lambda)),
\]
where
\[
\frac{\partial F}{\partial C}>0,
\qquad
\frac{\partial F}{\partial D}>0.
\]

If
\[
C(\lambda)
\]
continues increasing while
\[
D(\lambda)
\]
begins decreasing, then maximum aggregate strength need not coincide
with maximum useful registered output.

An interior optimum may occur.

This is not guaranteed by the finite theorem.

It is a structural possibility.

\subsection{Overcompression}

\begin{definition}[Overcompression]
Overcompression begins when additional relational loading that
previously improved coherent registration instead reduces the
independent contrast required by the surviving output mode.
\end{definition}

The relevant transition is therefore not simply
\[
\text{large norm}
\longrightarrow
\text{larger norm}.
\]

It is
\[
\text{organization that preserves distinction}
\]
to
\[
\text{organization that suppresses required distinction}.
\]

Overcompression is thus defined by a change in the role of added
constraint or loading.

\subsection{Normalized quadratic model}

The normalized Thalean quadratic gives a simple limiting example.

For
\[
\widehat Q=\frac{Q}{98},
\]
the common component is fixed:
\[
P_0\widehat Q^{\,k}x
=
P_0x.
\]

The centered component contracts:
\[
R(\widehat Q^{\,k}x)
\le
\left(\frac{9}{49}\right)^k
R(x).
\]

Therefore indefinite normalized \(\widehat Q\)-iteration drives every
centered component asymptotically toward zero.

If a useful outward registration requires nonzero centered structure,
then sufficiently deep \(\widehat Q\)-iteration necessarily lies on
the overcompressed side of that requirement.

This conclusion is conditional on the declared output requiring such
contrast.

\subsection{Clean release as a boundary regime}


The upgraded finite theorem does not establish a single monotone
sequence from native loading to contrast saturation.

Instead, two separately proved operations exist:
\[
\text{native selection}
\]
and
\[
\text{quadratic contrast concentration}.
\]

A larger composed dynamics might contain a sequence of the schematic
form
\[
\text{relational organization}
\]
\[
\downarrow
\]
\[
\text{native selection}
\]
\[
\downarrow
\]
\[
\text{coherent registered output}
\]
\[
\downarrow
\]
\[
\text{clean-release regime}
\]
\[
\downarrow
\]
\[
\text{contrast-concentrating stage}
\]
\[
\downarrow
\]
\[
\text{overcompression}
\]
\[
\downarrow
\]
\[
\text{contrast saturation}.
\]

This remains an internal hypothesis.

Both native loading and quadratic contrast concentration are now
theorem-level finite operations.

What remains unproved is the dynamical law that composes them into one
cycle and determines whether an intermediate optimum exists.

\subsection{A possible diagnostic}

For a declared register with
\[
C(h)\neq0,
\]
define the relative contrast ratio
\[
\rho(h)
=
\frac{R(h)}{C(h)}.
\]

A clean-release regime may be sought where:

\[
C(h)
\]
is large or increasing,

while
\[
\rho(h)
\]
remains above the minimum required for the declared output.

Overcompression begins when additional loading causes
\[
\rho(h)
\]
to fall below that functional regime.

This is a diagnostic proposal.

No universal critical value of
\[
\rho
\]
is asserted.

\subsection{Relation to projection-relative saturation}


Clean release, native future saturation, and contrast saturation are
not synonyms.

Native registered-future saturation occurs when
\[
|\mathcal C(p)|=1.
\]

Contrast saturation concerns suppression of a declared centered
contrast statistic.

Clean release concerns a possible useful outward-registration regime
in a larger composed dynamics.

The present theorem does not establish a universal ordering among
these three notions.

In particular, native loading cannot itself be called overcompression
merely because it adds constraint: under the canonical completion
register it increases, rather than decreases, centered contrast.

Any transition from future sharpening to contrast suppression requires
an additional certified composition law.

\subsection{No physical promotion}

The terms in this section are internal relational-mechanics terms.

They are not identified here with:

\begin{itemize}

\item sonoluminescence emission;

\item hydrodynamic collapse;

\item gravitational collapse;

\item black-hole formation;

\item Hawking radiation;

\item plasma coherence;

\item or any physical emission mechanism.

\end{itemize}

Those would require separate empirical and physical crosswalks.

\subsection{Boundary}


The exact theorem establishes:

\begin{itemize}

\item native relational loading with
\[
8\to2\to1;
\]

\item retained contextual distinction at registered-future
saturation;

\item the exact retained-sector response
\[
\Delta^{\mathsf T}\Delta=\frac34H;
\]

\item and independent quadratic centered contraction under
\(\widehat Q\).

\end{itemize}

Clean release and overcompression remain conditional concepts because
no theorem yet specifies the dynamics that compose native selection
with quadratic concentration into one loading cycle.

Neither concept is promoted to physical interpretation here.

\section{Empirical Falsification Boundary}

\subsection{Why an empirical boundary is needed}


The native finite Thalean relational-saturation theorem is now an
internal mathematical result.

The present paper does not claim empirical confirmation or physical
realization of that theorem.

The wider research program nevertheless contains a useful
methodological precedent: a native finite registration was converted
into a bounded scalar detector and subjected to a frozen
cross-instrument replication protocol.

That precedent shows how a future physical saturation test should be
designed.

It does not provide empirical evidence for the finite saturation
theorem itself.

\subsection{The frozen G60 C3 detector}

The WMAP replication \citep{cave2026wmap} reused the native
\(C_3\) registration inherited from the earlier Planck experiment.

For each accepted spherical placement and retained multipole band, the
detector evaluates three registered phase states under the native
\(C_3\) action.

Their harmonic response energies define a local three-phase profile.

That profile is mapped to a bounded scalar
\[
\Xi.
\]

The real-data pipeline is schematically
\[
T_{\ell_{\max}}
\longrightarrow
\nabla_{S^2}T_{\ell_{\max}}
\longrightarrow
\text{registered }C_3\text{ edge responses}
\longrightarrow
\text{three-phase energy profile}
\longrightarrow
\Xi.
\]

The exact graph registration, phase order, orientation, spherical
sampling geometry, edge-field construction, energy definition, and
\(\Xi\) evaluator were frozen before WMAP unblinding.

\subsection{Preregistered replication rule}

The frozen primary target was
\[
\ell_{\max}=64,
\qquad
\text{mean }\Xi,
\qquad
\text{upper tail}.
\]

Both WMAP V and W products were required independently to satisfy
\[
p_{\mathrm{two}}\le0.05
\]
against both of their frozen null families:

\begin{enumerate}[label=(\roman*)]

\item Gaussian nulls;

\item amplitude-preserving phase-surrogate nulls.

\end{enumerate}

Thus all four required product-null comparisons had to pass.

No combined probability was permitted.

No neighboring scale, alternate reduction, secondary product, or
post-hoc statistic could rescue a failed primary comparison.

\subsection{Observed result}

The four frozen primary comparisons were

\[
p_{\mathrm{V,Gaussian}}
=
0.0239760,
\]

\[
p_{\mathrm{V,phase}}
=
0.00799201,
\]

\[
p_{\mathrm{W,Gaussian}}
=
0.0539461,
\]

and

\[
p_{\mathrm{W,phase}}
=
0.0339660.
\]

All four real WMAP statistics lay in the preregistered upper-tail
direction.

Three of the four comparisons crossed
\[
p\le0.05.
\]

The WMAP W Gaussian comparison did not.

Therefore the frozen decision rule yields
\[
\boxed{\text{NO-GO REPLICATION}}.
\]

No post-hoc rescue is applied.

\subsection{What the WMAP result does and does not show}


The result preserves one descriptive observation:

\begin{quote}
The predicted upper-tail direction appeared in all four required WMAP
comparisons.
\end{quote}

But the bounded cross-instrument replication claim was not earned.

The finite mathematical construction is unchanged by that outcome.

The WMAP paper therefore provides a methodological lesson, not
empirical support for native relational saturation, quadratic contrast
concentration, or any physical saturation bridge.

\subsection{No retroactive saturation detector}

The scalar
\[
\Xi
\]
was not preregistered as a relational-saturation observable.

It measured a frozen \(C_3\) phase-structure response.

Therefore Paper 14 does not reinterpret
\[
\Xi
\]
after the fact as:

\begin{itemize}

\item a saturation measure;

\item a gravity measure;

\item a black-hole proxy;

\item an information-density measure;

\item or a contrast-concentration statistic.

\end{itemize}

Doing so would define a new experiment.

\subsection{Requirements for a future saturation test}


A future empirical test must first declare which finite theorem-level
operation is being mapped into physics.

At minimum it should freeze:

\begin{enumerate}[label=(\roman*)]

\item the physical or complete mathematical state family;

\item the declared projection or registration;

\item the proposed physical analogue of native selection, if any;

\item the proposed contrast-concentrating operation, if any;

\item the aggregate observable;

\item the registered contrast observable;

\item the expected relation among those quantities;

\item the null model or comparison family;

\item the predicted direction;

\item the decision threshold;

\item and the rule for classifying the outcome.

\end{enumerate}

Native relational loading and quadratic contrast concentration may not
be treated as one undifferentiated loading variable.

A proposed physical cycle containing both must specify their coupling
before observing the data.

\subsection{A minimal structural observable}


Suppose a physical bridge supplies a declared register
\[
h(\lambda)
\]
indexed by a physical control parameter
\[
\lambda.
\]

Define
\[
C(\lambda)=\|P_0h(\lambda)\|_2
\]
and
\[
\zeta(\lambda)
=
\frac{\|P_\perp h(\lambda)\|_2^2}
{\|h(\lambda)\|_2^2}.
\]

A direct test of a \emph{contrast-concentrating} physical regime might
predict that
\[
C(\lambda)
\]
remains large or increases while
\[
\zeta(\lambda)
\]
decreases.

That is not the theorem-level signature of native relational loading.

Under the canonical finite completion register, native loading
strictly increases \(\zeta\).

Therefore any physical experiment must state whether it is testing:

\[
\text{native future selection},
\]
\[
\text{contrast concentration},
\]
or
\[
\text{an explicitly proposed composition of the two}.
\]

The existence of an observed pattern would test the declared bridge.

It would not, by itself, establish its physical cause.

\subsection{Falsification logic}


A future physical saturation hypothesis should count against itself if
one or more of the following occurs under a valid frozen protocol:

\begin{enumerate}[label=(\roman*)]

\item the declared native-selection analogue fails to produce its
predicted change in future or state multiplicity;

\item the declared contrast-concentrating analogue fails to produce
its predicted contrast behavior;

\item a proposed coupling between the two is not observed;

\item the observed relation is reproduced equally by declared null
structures lacking the relevant Thalean organization;

\item the signal depends on post-hoc changes to registration, loading,
scale, coupling law, or threshold;

\item or the proposed effect fails independent replication.

\end{enumerate}

The purpose is to keep the physical extension falsifiable rather than
interpretively elastic.

\subsection{Boundary}


The exact empirical status carried into Paper 14 is:

\[
\text{native finite detector methodology}
\quad
\text{demonstrated},
\]

\[
\text{WMAP directional concordance}
\quad
\text{observed},
\]

\[
\text{WMAP preregistered replication}
\quad
\text{NO-GO},
\]

\[
\text{physical relational saturation}
\quad
\text{untested}.
\]

No empirical result currently promotes the native finite
relational-saturation theorem to a physical claim.

\section{Physical Bridge Boundary}

The theorem established in this paper is finite relational
mathematics.

Native relational saturation is now an internal theorem of the
registered finite mechanics.

Quadratic contrast concentration is a second finite theorem-level
operation.

Physical interpretation belongs to a separate layer.

No result in the preceding sections is promoted to physical gravity,
black-hole physics, spacetime, or Planck-scale information theory
without an additional bridge theorem.

\subsection{Three layers}


The epistemic structure of Paper 14 is now
\[
\text{finite relational theorem}
\longrightarrow
\text{internal conditional extensions}
\longrightarrow
\text{physical bridge hypotheses}.
\]

The first layer contains native relational loading, registered-future
saturation, retained projection-relative distinction, and quadratic
contrast concentration.

The second contains proposed compositions such as clean release and
overcompression, together with stronger global transport
constructions.

The third requires an explicit mapping from finite theorem objects to
physical observables.

No physical noun promotes an internal definition.

Thus the phrases
\[
\text{relational saturation},
\qquad
\text{registration boundary},
\qquad
\text{clean release},
\qquad
\text{overcompression}
\]
remain internal terms unless a physical bridge is separately proved.

\subsection{Finite distinction support}

One possible physical extension begins from the hypothesis that an
independent physically realizable distinction-event requires finite
support.

Let
\[
r_E(\mathcal R)
\]
denote the number or rank of independent actionable distinction-events
realized in a physical region
\[
\mathcal R.
\]

A physical finite-density hypothesis might take the form
\[
r_E(\mathcal R)
\le
\rho_{E,\max}V_4(\mathcal R),
\]
where
\[
V_4(\mathcal R)
\]
is an independently defined physical four-volume.

A stronger dimensional conjecture might propose
\[
\rho_{E,\max}
\sim
\ell_P^{-4}.
\]

Equivalently, one might conjecture a minimum physical four-support of
order
\[
\ell_P^4
\]
per independent physically realizable distinction-event.

None of these statements is proved in Paper 14.

In particular, the paper does not assert
\[
1\text{ bit}=\ell_P^4.
\]

The relevant conjectural quantity is independent distinction density,
not arbitrary binary labeling.

\subsection{Event capacity is not static storage}

A four-volume density, if physically meaningful, cannot be identified
without argument with the static information capacity of a stationary
object.

A persistent physical system may occupy increasing four-volume as
duration accumulates while its instantaneous state capacity remains
unchanged.

Therefore any future
\[
\ell_P^4
\]
bridge must distinguish at least:

\begin{enumerate}[label=(\roman*)]

\item state capacity;

\item event or innovation capacity;

\item history capacity;

\item boundary-transmissible registration capacity.

\end{enumerate}

Paper 14 does not collapse these categories.

\subsection{Thalean gravity remains a bridge hypothesis}


The upgraded finite architecture separates two operations that earlier
versions left partially conflated:

\[
\text{native future selection}
\]
and
\[
\text{quadratic contrast concentration}.
\]

A future physical gravity hypothesis would therefore have to identify:

\begin{enumerate}[label=(\roman*)]

\item a physical state space or observable family;

\item a map from that physical structure to the finite relational
register;

\item a physical analogue of native future selection;

\item a physical analogue of contrast concentration, if required;

\item an explicit coupling law if both participate in one physical
process;

\item a derivation of the relevant aggregate and registered contrast;

\item and a demonstrated relation to gravitational observables.

\end{enumerate}

No such derivation is supplied here.

Therefore
\[
\text{finite relational saturation}
\not\Rightarrow
\text{physical gravity}.
\]

\subsection{No nonlinear graph deformation is introduced}

The finite mechanics used in Paper 14 is linear.

The graph and the quadratic operator are not deformed as a function
of loading.

The theorem does not propose a state-dependent metric
\[
G_{ij}(x),
\]
a nonlinear adjacency law, or a graph that physically bends under
source strength.

The exact operator remains
\[
Q=A^3+2A^2+2I.
\]

Any future physical gravity bridge must therefore respect the linear
finite substrate unless an independently justified extension is
introduced.

\subsection{Long range is not a separate selector}

Paper 14 also does not introduce a universal long-range gravitational
selector.

The internal QR architecture instead seeks long-range structure
through lawful accumulation over an admitted summation domain.

Schematically,
\[
\text{local admissible contributions}
\longrightarrow
\text{global compatibility}
\longrightarrow
\text{lawful sum}.
\]

A physical bridge would need to show how such a finite sum produces
the appropriate macroscopic gravitational behavior.

That bridge has not been derived.

\subsection{Black-hole interpretation}

Relational saturation may motivate a later black-hole hypothesis.

The internal mechanism would be:

\[
\text{finite complete distinction}
\]
together with
\[
\text{increasing aggregate dominance}
\]
and
\[
\text{decreasing outward individuation}.
\]

This provides an internal finite alternative to requiring unbounded
independent distinction density.

It does not prove that an astrophysical black hole realizes this
mechanism.

Accordingly, Paper 14 does not identify:

\begin{itemize}

\item contrast saturation with gravitational collapse;

\item a projection fibre with a black-hole interior;

\item a registration boundary with an event horizon;

\item common-mode dominance with mass density;

\item or finite relational degeneracy with a physical singularity.

\end{itemize}

These remain possible bridge hypotheses only.

\subsection{Horizon-like registration}

Within the internal mechanics one may define a horizon-like condition
relative to outward registration.

Let
\[
\Pi_{\mathrm{out}}
\]
denote a declared family of outward registrations.

A finite internal system is registration-saturated when additional
complete distinctions fail to produce additional independently
distinguishable outward registrations.

This can be stated schematically as
\[
r_{\mathrm{int}}
>
r_{\mathrm{out}}.
\]

That statement is internal.

A physical event horizon is a causal object in spacetime.

No theorem in Paper 14 identifies the two.

\subsection{No singularity theorem}

The finite contrast-concentration theorem contains no divergent
quantity.

Instead,
\[
\widehat Q^{\,k}
\]
remains finite for every
\[
k.
\]

Its centered component contracts, and its operator limit is the
finite rank-one projection
\[
P_0.
\]

This gives a finite algebraic alternative to divergence inside the
model.

It does not establish singularity avoidance in general relativity.

A physical singularity claim would require an independently derived
spacetime bridge and the appropriate causal, metric, and geodesic
analysis.

Therefore
\[
\text{finite saturation}
\not\Rightarrow
\text{physical singularity resolution}.
\]

\subsection{Black-hole entropy}

Paper 14 does not derive the Bekenstein-Hawking relation
\[
S_{\mathrm{BH}}
=
\frac{k_BA}{4\ell_P^2}.
\]

In particular, the conjectured four-support
\[
\ell_P^4
\]
for independent event distinction is not the same mathematical
quantity as an area-scaled black-hole entropy.

Any future relation between these ideas would have to explain:

\begin{enumerate}[label=(\roman*)]

\item why a four-dimensional event-capacity law produces an
area-scaled boundary law;

\item why the relevant coefficient is
\[
\frac14;
\]

\item which physical distinctions are being counted;

\item how the boundary transmits or registers them.

\end{enumerate}

None of those derivations is supplied here.

\subsection{No Hawking-radiation claim}

The clean-release language introduced in Section 11 is not identified
with Hawking radiation.

Clean release is an internal statement about coherent registered
output and surviving contrast.

Hawking radiation is a physical quantum-field-theoretic phenomenon
requiring a separate physical derivation.

No bridge between them is asserted.

\subsection{Cosmic microwave background}

A later QR interpretation may ask whether cosmological transparency
can be understood as an opening of a larger registration or
summation domain.

Paper 14 does not establish that interpretation.

The existing WMAP experiment tested a frozen native \(C_3\)
phase-structure detector and returned a formal NO-GO replication.

It therefore cannot be used here as evidence for a physical
saturation mechanism.

A future CMB saturation experiment would require a new frozen
observable and a new decision contract.

\subsection{Thalean time remains intrinsic finite time}

The established finite theory defines
\[
\tau_{\mathrm T}
=
\operatorname{wind}_{r_1}
\]
as the primitive additive oriented winding of the unique
\(G_{30}\)-visible registered-history cycle.

Paper 14 does not identify
\[
\tau_{\mathrm T}
\]
with:

\begin{itemize}

\item seconds;

\item proper time;

\item physical duration;

\item a Lorentzian temporal coordinate;

\item one quantum of action;

\item \(h\) or \(\hbar\);

\item or a spacetime interval.

\end{itemize}

No physical calibration is supplied by the saturation theorem.

\subsection{Physical falsifiability requirement}


A future physical bridge must make more than a qualitative analogy.

It must identify observables and predict a measurable relation that
could fail.

Because the finite theory now distinguishes native selection from
contrast concentration, the physical bridge must do the same.

A contrast-concentration prediction might have the form
\[
C_{\mathrm{phys}}(\lambda)
\quad
\text{remains large or increases}
\]
while
\[
\zeta_{\mathrm{phys}}(\lambda)
\quad
\text{decreases}.
\]

A native-selection prediction would require a separately defined
physical analogue of decreasing lawful future freedom.

A proposed physical cycle containing both must preregister the
coupling between them.

Without those mappings, the finite theorem remains mathematics rather
than physics.

\subsection{Explicit nonclaims}

Paper 14 does not establish:

\begin{enumerate}[label=(\roman*)]

\item physical information density;

\item one bit per Planck four-volume;

\item physical gravity;

\item Newtonian gravity;

\item general relativity;

\item the equivalence principle;

\item physical spacetime;

\item black-hole interiors;

\item physical event horizons;

\item singularity avoidance;

\item black-hole entropy;

\item holography;

\item Hawking radiation;

\item a physical CMB saturation signal;

\item or a Standard Model embedding.

\end{enumerate}

\subsection{Bridge keeper}


The relationship among the layers can be summarized as follows:

\begin{quote}
Finite relational mathematics can support internal mechanical
extensions.  Internal mechanics can motivate physical hypotheses.
Physical claims do not follow without an explicit physical bridge.
\end{quote}

Paper 14 now establishes both a native finite relational-saturation
theorem and a separate quadratic contrast-concentration theorem.

Clean-release composition, stronger global inter-face transport, and
all physical interpretations remain additional hypotheses.

Whether nature realizes any analogue of this finite architecture
remains an open physical question.

\section{Discussion}

\subsection{The theorem in one sentence}

The finite result of Paper 14 can now be stated as follows:

\begin{quote}
Admissible native transport reduces lawful future freedom from
\(8\) to \(2\) to \(1\), while contextual distinction remains
retained; the resulting designation response selects exactly the
ten-dimensional five-address-hidden sector, and the Thalean
quadratic independently contracts centered registered contrast
without finite-stage erasure.
\end{quote}

The important word is \emph{independently}.

Native relational loading and quadratic contrast concentration are
not two descriptions of the same operation.

They are complementary finite mechanisms.

\subsection{Selection, retention, and concentration}

The mature architecture separates three operations.

\paragraph{Selection.}
For a live prefix \(p\), let
\[
\mathcal C(p)
\]
be the lawful completion family.

Native loading is ordered by
\[
p\preceq p'
\quad\Longrightarrow\quad
\mathcal C(p')\subseteq\mathcal C(p).
\]

For the certified carrier,
\[
8\longrightarrow2\longrightarrow1.
\]

Thus admitted transport progressively selects a future.

\paragraph{Retention.}
Future uniqueness does not collapse the complete relational state.

At the two-edge threshold there remain \(120\) distinct contexts
with \(120\) distinct completed words.

The native designation response
\[
\Delta
\]
acts entirely in the retained five-address-hidden sector:
\[
\Delta^{\mathsf T}\Delta
=
\Delta\Delta^{\mathsf T}
=
\frac34H,
\]
where
\[
\operatorname{rank}(H)=10.
\]

The declared five-address readout satisfies
\[
\rho\Delta=0.
\]

Thus the retained response is nontrivial internally while invisible
to the coarse readout.

\paragraph{Concentration.}
The normalized quadratic
\[
\widehat Q=\frac1{98}Q
\]
fixes common mode and contracts every centered eigendirection.

Its sharp centered factor is
\[
\frac9{49}.
\]

Consequently,
\[
R(\widehat Q^{\,k}x)
\le
\left(\frac9{49}\right)^k R(x),
\]
while every finite iterate remains invertible.

Selection, retention, and concentration are therefore distinct
finite operations with exact algebraic relations.

\subsection{Native loading sharpens rather than smooths}

The canonical completion-occupancy register provides a decisive
test of the relation between native loading and Paper 14 contrast.

Its exact contrast values are
\[
\zeta:
\frac{31}{55}
\longrightarrow
\frac{41}{65}
\longrightarrow
\frac{7}{10}.
\]

Likewise,
\[
\rho^2:
\frac{31}{24}
\longrightarrow
\frac{41}{24}
\longrightarrow
\frac73.
\]

The common component remains fixed:
\[
C^2=\frac{12}{5}.
\]

All \(60\) strict \(8\to2\) extensions increase both contrast
statistics.

All \(120\) strict \(2\to1\) extensions do the same.

There are no exceptions.

Thus the native loading process is contrast-antimonotone relative to
the normalized quadratic.

The correct interpretation is:

\[
\text{native loading}
\longrightarrow
\text{future sharpening},
\]
whereas
\[
\widehat Q
\longrightarrow
\text{common-mode concentration}.
\]

This distinction prevents a false identification that earlier
versions of the saturation conjecture left open.

\subsection{The retained-sector identity}

The strongest native-to-register crosswalk is not
\[
T\mapsto\widehat Q.
\]

It is
\[
\Delta^{\mathsf T}\Delta=\frac34H.
\]

This identity says that the quadratic response of context-activated
native designation is exactly the projector onto the retained
ten-dimensional relational sector, up to the scalar \(3/4\).

Because also
\[
\Delta\Delta^{\mathsf T}=\frac34H,
\]
the normalized operator
\[
\sqrt{\frac43}\,\Delta
\]
is an isometry on the retained sector.

Thus native contextual transport carries an exact internally
nontrivial response even though the five-address readout sees
\[
\rho\Delta=0.
\]

This is a finite realization of retained distinction behind
registration.

\subsection{Relation to the Thalean response algebra}

The native loading Gram belongs to the same exact centered
transport-response algebra as the centered Thalean quadratic.

If \(R_0,R_1,R_2\) are the three certified centered response classes,
then
\[
\Delta^{\mathsf T}\Delta
=
\frac1{48}R_0
-
\frac1{16}R_1
+
\frac1{32}R_2.
\]

The centered Thalean quadratic satisfies
\[
PQP
=
\frac{139}{90}R_0
-
\frac{28}{15}R_1
+
\frac{19}{60}R_2.
\]

The operators are not proportional.

Their relationship is therefore algebraic rather than identificatory:
native loading and quadratic concentration inhabit the same certified
response algebra while selecting different directions within it.

\subsection{Two distinct saturation notions}

The upgraded theory contains two theorem-level saturation notions.

\paragraph{Registered-future saturation.}
A native prefix is saturated when
\[
|\mathcal C(p)|=1.
\]

This means lawful future freedom has become unique.

It does not mean complete relational distinction has disappeared.

\paragraph{Contrast saturation.}
A registered sequence is contrast-saturating when its declared
contrast approaches a stable minimum while complete state remains
nontrivial.

Normalized quadratic aggregation supplies the exact example
\[
\zeta(\widehat Q^{\,k}x)\longrightarrow0
\]
whenever the input has a nonzero common component.

Thus
\[
\text{registered-future saturation}
\neq
\text{contrast saturation}.
\]

The first is selection of a unique lawful future.

The second is suppression of centered registered contrast.

\subsection{No transfer into the common mode}

Contrast concentration under \(\widehat Q\) must still not be read as
transfer from the centered sector into common mode.

The common and centered subspaces are invariant.

If
\[
P_0x=0,
\]
then
\[
P_0\widehat Q^{\,k}x=0
\]
for every \(k\).

A purely centered input remains centered and contracts toward zero.

The concentration theorem is therefore a statement about differential
spectral persistence.

This remains distinct from the native carrier mechanism, where
transport creates contextual structure and designation becomes
operational after the first admitted step.

\subsection{Invertible stages and projection loss}

Every finite power of \(\widehat Q\) remains invertible.

Yet
\[
\widehat Q^{\,k}\longrightarrow P_0,
\]
and
\[
\operatorname{rank}(P_0)=1.
\]

Thus exact finite-stage distinction can remain recoverable while a
declared contrast readout becomes arbitrarily insensitive to it.

The five-address theorem supplies a stronger exact form of the same
general principle.

The hidden projector \(H\) has rank ten.

States differing entirely within that sector can remain distinct in
the fifteen-state register while having identical five-address
readout.

Projection loss and finite-stage invertibility are therefore two
different mechanisms by which visible individuation can weaken
without complete-state merger.

\subsection{Global sigma is a separate frontier}

The unresolved inter-face sigma remains mathematically interesting,
but it is no longer a missing premise of the finite saturation
theorem.

The native registered carrier used in Sections 8 and 14 is already
an admitted finite transport system with exact successor and
completion laws.

The sigma problem instead asks whether a stronger inter-face seam
transport can be made globally compatible under its declared
four-step closure conditions.

Therefore
\[
\text{sigma open}
\]
does not imply
\[
\text{native saturation open}.
\]

The current sigma frontier remains:

\[
2816
\]
neighbor-frame-\(E\)-closed local options,

\[
376
\]
completed solver iterations,

\[
175069
\]
unique nogoods,

and a best observed partial closure of
\[
32
\]
flags, without either a complete global \(r_3\) or a proof of
infeasibility.

That is a separate global-compatibility problem.

\subsection{Long range as lawful accumulation}

The finite theory still supports the architectural principle that
long-range registered structure should arise from lawful finite
accumulation rather than from an independently inserted long-range
selector.

The exact examples remain:

\[
x\longmapsto Qx
\]
on the sector register,

\[
d\delta=\omega
\]
on the binary closure complex,

\[
d\eta=\Omega
\]
on its integer refinement,

and
\[
\tau_{\mathrm T}(\gamma)
=
\operatorname{wind}_{r_1}(\gamma)
\]
on registered-history cycles.

Their domains and coefficient systems differ.

No theorem identifies them as one universal sum.

\subsection{Clean release and overcompression remain conditional}

The new native loading theorem changes the status of the clean-release
picture.

It is no longer correct to assume that increasing native loading
itself drives centered contrast downward.

For the canonical native completion register, it does the opposite.

A larger cycle containing both selection and concentration might still
possess an intermediate clean-release or overcompression regime, but
such a composition requires an explicit law specifying when and how
the native selection dynamics couple to a contrast-concentrating
operation.

No such cycle is proved in this paper.

Clean release and overcompression therefore remain internal
conditional concepts.

\subsection{Finite capacity remains a separate hypothesis}

The native loading theorem supplies a finite completion family and an
exact finite distinction count.

It does not prove a universal bound on all possible real-valued
relational states.

The fifteen-state sector space remains a real vector space.

Therefore finite graph dimension must not be silently identified with
finite physical information capacity.

Any stronger capacity statement requires a declared notion of
independent distinction, such as finite cardinality, linear rank,
finite-resolution distinguishability, or another explicitly defined
independence structure.

\subsection{No singularity is required internally}

Nothing in the finite theorem requires a divergent variable.

Native future freedom reaches
\[
|\mathcal C|=1
\]
at a finite stage.

The retained response remains finite:
\[
\Delta^{\mathsf T}\Delta=\frac34H.
\]

The normalized quadratic remains finite at every iterate and converges
to the finite projector
\[
P_0.
\]

Thus the internal mathematics reaches saturation through finite
selection, retained distinction, and projection-relative concentration
rather than through divergence.

This is not a theorem about singularities in general relativity.

\subsection{Methodological posture}

The upgraded result preserves the methodological rule:

\begin{quote}
Do not protect the idea. Make the idea protect itself.
\end{quote}

The canonical native loading test could have confirmed the earlier
expectation that loading itself decreases contrast.

It did not.

Every one of the \(180\) strict extensions moved contrast in the
opposite direction.

The correct response is not to change the readout after seeing the
result.

It is to change the theory.

That negative test is therefore part of the theorem architecture:
native selection and quadratic concentration must remain distinct.

\subsection{Current theorem frontier}

The finite internal status is now:

\[
\text{quadratic contrast concentration}
\quad
\text{proved},
\]

\[
\text{finite-stage quadratic non-erasure}
\quad
\text{proved},
\]

\[
\text{native relational loading }8\to2\to1
\quad
\text{proved},
\]

\[
\text{registered-future saturation}
\quad
\text{proved},
\]

\[
\text{five-address projection-relative saturation}
\quad
\text{proved},
\]

\[
\Delta^{\mathsf T}\Delta=\frac34H
\quad
\text{proved},
\]

\[
\text{native loading contrast antimonotonicity}
\quad
\text{proved},
\]

\[
\text{global inter-face sigma}
\quad
\text{open},
\]

\[
\text{clean-release composition law}
\quad
\text{open},
\]

\[
\text{physical saturation bridge}
\quad
\text{open}.
\]

The native finite relational-saturation theorem is therefore closed.

The remaining frontier concerns stronger composition laws and physical
embedding, not the existence of the finite saturation mechanism.

\section{Conclusion}

Paper 14 now establishes a native finite theorem of Thalean relational
saturation.

The result does not identify native loading with iteration of the
Thalean quadratic.

Instead, it separates three exact finite operations:
\[
\text{selection},
\qquad
\text{retention},
\qquad
\text{concentration}.
\]

Native selection is supplied by the certified registered carrier.

For a live prefix \(p\), the lawful completion family contracts
monotonically:
\[
8\longrightarrow2\longrightarrow1.
\]

Registered-future saturation occurs when
\[
|\mathcal C(p)|=1.
\]

Future uniqueness does not erase complete relational distinction.

At the two-edge threshold there remain \(120\) distinct contexts with
\(120\) distinct registered completions, and the terminal completion
register retains two native holonomy designations over each closing
chamber.

Transport also creates the context in which native designation becomes
operational.

If \(\Delta\) is the induced context-activated designation response
and \(H\) the ten-dimensional five-address-hidden projector, then
\[
\Delta^{\mathsf T}\Delta
=
\Delta\Delta^{\mathsf T}
=
\frac34H,
\]
with
\[
\rho\Delta=0.
\]

Thus the native response is exactly supported on a retained relational
sector that is invisible to the declared coarse readout.

Native loading does not itself perform common-mode contrast
concentration.

For the canonical completion-occupancy register,
\[
\zeta:
\frac{31}{55}
\longrightarrow
\frac{41}{65}
\longrightarrow
\frac{7}{10},
\]
and
\[
\rho^2:
\frac{31}{24}
\longrightarrow
\frac{41}{24}
\longrightarrow
\frac73.
\]

Every one of the \(180\) strict native loading extensions increases
these contrast statistics.

Native loading therefore sharpens the selected future.

The normalized Thalean quadratic supplies the distinct concentrating
operation:
\[
\widehat Q=\frac1{98}Q,
\qquad
Q=MM^{\mathsf T}=A^3+2A^2+2I.
\]

Its common mode is fixed, every centered mode contracts, and
\[
R(\widehat Q^{\,k}x)
\le
\left(\frac9{49}\right)^kR(x).
\]

Every finite iterate remains invertible, even though
\[
\widehat Q^{\,k}\longrightarrow P_0.
\]

The native loading response and the centered quadratic belong to the
same certified transport-response algebra, but they are not the same
operator and do not perform the same finite action.

\subsection{What is theorem}

The following statements are established internally:

\begin{enumerate}[label=(\roman*)]

\item native relational loading is monotone contraction of the lawful
completion family;

\item the exact loading profile is
\[
8\to2\to1;
\]

\item registered-future saturation occurs at the two-edge threshold;

\item unique future does not imply unique complete context;

\item context activates native \(b/ab\) designation after one admitted
carrier step;

\item the induced designation response has rank ten;

\item
\[
\Delta^{\mathsf T}\Delta
=
\Delta\Delta^{\mathsf T}
=
\frac34H;
\]

\item the retained response is invisible to the five-address readout;

\item the canonical completion register strictly increases Paper 14
contrast along all \(180\) strict loading extensions;

\item native future saturation and contrast saturation are distinct;

\item the normalized Thalean quadratic contracts centered contrast by
the sharp factor \(9/49\);

\item every finite quadratic iterate remains invertible;

\item the native loading response and the centered quadratic inhabit
the same exact transport-response algebra.

\end{enumerate}

Together these statements close the native finite
relational-saturation theorem.

\subsection{What remains conditional internally}

The following are not promoted to theorem:

\begin{itemize}

\item a unique clean-release optimum;

\item a universal overcompression threshold;

\item a dynamical cycle specifying when native selection is followed
by or coupled to quadratic concentration;

\item a globally closed inter-face sigma;

\item a universal finite information-capacity law.

\end{itemize}

The open sigma problem belongs to a stronger global inter-face
transport construction.

It is not a missing premise of the finite saturation theorem.

\subsection{The theorem ladder}

The closed finite chain is now

\[
\text{native registered carrier}
\]
\[
\downarrow
\]
\[
\text{context activation}
\]
\[
\downarrow
\]
\[
8\to2\to1
\]
\[
\downarrow
\]
\[
\text{registered-future saturation}
\]
\[
\downarrow
\]
\[
\Delta^{\mathsf T}\Delta=\frac34H
\]
\[
\downarrow
\]
\[
\text{retained projection-relative distinction}.
\]

Separately,
\[
\text{Thalean incidence algebra}
\]
\[
\downarrow
\]
\[
\widehat Q
\]
\[
\downarrow
\]
\[
\text{centered contrast contraction}
\]
\[
\downarrow
\]
\[
\text{finite-stage non-erasure}.
\]

The two chains meet in the common certified centered
transport-response algebra.

They are not collapsed into one operation.

\subsection{Physical boundary}

Nothing in the finite theorem identifies relational saturation with:

\begin{itemize}

\item physical gravity;

\item Newtonian or relativistic gravitational dynamics;

\item black-hole collapse;

\item an event horizon;

\item singularity avoidance;

\item black-hole entropy;

\item Hawking radiation;

\item Planck-scale information density;

\item physical proper time;

\item the cosmic microwave background;

\item or a Standard Model interaction.

\end{itemize}

Those require separate physical bridge theorems and empirical tests.

\subsection{Final statement}

The central finite lesson is:

\begin{quote}
Relational saturation is reached when admissible future freedom
becomes unique while complete relational distinction remains retained
behind registration.
\end{quote}

The exact native signature is
\[
|\mathcal C|:
8\longrightarrow2\longrightarrow1,
\]
together with
\[
\Delta^{\mathsf T}\Delta=\frac34H,
\qquad
\rho\Delta=0.
\]

The exact quadratic signature is
\[
R(\widehat Q^{\,k}x)
\le
\left(\frac9{49}\right)^kR(x)
\]
with finite-stage invertibility.

The mature theory therefore does not say that greater native loading
automatically produces lower contrast.

It says something more structured:

\begin{quote}
Selection sharpens a lawful future.
Context retains distinctions hidden by coarse registration.
Quadratic aggregation can independently suppress registered contrast
without erasing the retained state.
\end{quote}

Whether nature realizes any physical analogue of this finite
architecture remains an open question.

\section{Theorem and Provenance Ledger}

This section records the exact status of the structures used in
Paper 14.

The purpose is to keep imported theorem-level facts, new results,
internal conjectures, computational frontiers, and physical bridge
hypotheses distinct.

\subsection{Imported theorem-level structure}

The following objects are imported from the established finite
AT4val[60,6] theorem chain.

\begin{center}
\begin{tabular}{p{0.34\linewidth}p{0.56\linewidth}}
\toprule
Object or statement & Status \\
\midrule

\(
G_{15}\cong L(\mathrm{Petersen})
\)
&
proved finite graph identification
\\

\(
M\in\{0,1\}^{15\times30}
\)
&
proved transport-induced sector-edge incidence matrix
\\

row weight \(14\), column weight \(7\)
&
proved incidence census
\\

\(
Q=MM^{\mathsf T}
\)
&
proved sector-overlap Gram identity
\\

\(
Q=A^3+2A^2+2I
\)
&
proved adjacency-algebra identity
\\

\(
\operatorname{spec}(Q)
=
\{98^{(1)},18^{(5)},3^{(4)},2^{(5)}\}
\)
&
proved spectral decomposition
\\

signed
\(
G_{30}\to G_{15}
\)
lift
&
proved finite covering structure
\\

nontrivial binary switching class
&
proved lifted holonomy not determined by \(Q\)
\\

projected return before complete return
&
proved in multiple declared finite constructions
\\

registered history
&
proved retained finite distinction relative to declared projections
\\

\(
\tau_{\mathrm T}
=
\operatorname{wind}_{r_1}
\)
&
proved intrinsic additive oriented history coordinate
\\

\(
R_n(h)
=
\|h-\overline h\,\mathbf 1_n\|_2
\)
&
proved centered Euclidean invariant on a declared finite register
\\

\bottomrule
\end{tabular}
\end{center}

\subsection{New theorem-level results of Paper 14}

Paper 14 introduces the canonical normalization

\[
\widehat Q
=
\frac{1}{98}Q.
\]

The following results are proved in this paper.

\begin{center}
\begin{tabular}{p{0.42\linewidth}p{0.48\linewidth}}
\toprule
Statement & Status \\
\midrule

\(
\widehat Q\mathbf 1
=
\mathbf 1
\)
&
proved
\\

common-mode eigenvalue
\(
1
\)
&
proved
\\

centered eigenvalues
\(
9/49,\ 3/98,\ 1/49
\)
&
proved
\\

centered spectral radius
\[
\rho_{\mathrm{cent}}
=
\frac{9}{49}
\]
&
proved
\\

\[
R(\widehat Q^{\,k}x)
\le
\left(\frac{9}{49}\right)^kR(x)
\]
&
proved
\\

\[
\widehat Q^{\,k}
\longrightarrow
P_0
\]
&
proved in Euclidean operator norm
\\

\[
\left\|
\widehat Q^{\,k}-P_0
\right\|_{2\to2}
=
\left(\frac{9}{49}\right)^k
\]
&
proved and sharp
\\

\(
\det\widehat Q\neq0
\)
&
proved
\\

finite-stage invertibility of
\(
\widehat Q^{\,k}
\)
&
proved for every finite \(k\)
\\

rank-one asymptotic limit
&
proved because
\(
\operatorname{rank}(P_0)=1
\)
\\

finite-resolution contrast loss
&
proved relative to any fixed positive threshold
\\

\bottomrule
\end{tabular}
\end{center}

\subsection{Exact theorem keeper}

The new theorem package may be summarized as:

\begin{quote}
The normalized Thalean quadratic is an invertible finite linear
contrast concentrator whose iterates converge sharply to the
common-mode projection.
\end{quote}

Equivalently:

\begin{quote}
Relative registered contrast may approach zero while exact
finite-stage linear distinction remains preserved.
\end{quote}

\subsection{Defined but not promoted to theorem}

The following terms are defined in Paper 14 but are not themselves
proved native dynamical laws:

\begin{itemize}

\item projection-relative saturation;

\item contrast saturation;

\item effective registration degeneracy;

\item clean release;

\item overcompression;

\item inter-face sigma;

\item admitted summation domain.

\end{itemize}

These definitions organize the conjectural mechanics around the exact
linear theorem.

\subsection{Thalean Relational Saturation Conjecture}

The central conjectural extension is:

\begin{quote}
Given a globally compatible finite linear admissible transport
structure, a declared projection or registration, a finite additive
aggregation law, and finite independent distinction capacity,
increasing relational loading can strengthen the coherent aggregate
of bounded admissible contributions while reducing their independent
registered contrast.

Beyond a critical regime, additional loading is expressed through
increasing dependency, common-mode concentration, and
projection-relative degeneracy rather than unbounded growth of
independent local distinction.

Complete distinctions may remain preserved by the finite transport
while the declared registration no longer resolves them independently.
\end{quote}

This statement is not yet a theorem of the native higher mechanics.

\subsection{Missing native crosswalk}

The principal missing mathematical arrow is

\[
\text{native admitted transport/loading}
\longrightarrow
\text{certified contrast-concentrating register}.
\]

A strong theorem would require a diagram of the form

\[
\begin{array}{ccc}
X_k & \xrightarrow{T} & X_{k+1} \\
\downarrow \rho & & \downarrow \rho \\
H_k & \xrightarrow{L} & H_{k+1},
\end{array}
\]

with

\[
\rho\circ T
=
L\circ\rho,
\]

where:

\begin{enumerate}[label=(\roman*)]

\item \(T\) is a native admitted transformation;

\item \(\rho\) is a certified registered readout;

\item \(L\) is contrast-concentrating;

\item retained complete distinction remains nontrivial upstairs.

\end{enumerate}

Paper 14 proves the required behavior for the candidate operator

\[
L=\widehat Q.
\]

The native \(T\) and \(\rho\) crosswalk remain open.

\subsection{Current inter-face sigma frontier}

The completed Project 41 Mac search used the canonically completed,
neighbor-frame-\(E\)-closed local seam vocabulary containing

\[
2816
\]

local options.

The completed search reached:

\begin{itemize}

\item \(376\) total solver iterations;

\item \(175069\) unique four-step closure nogoods;

\item repeated best candidates with
\[
1888
\]
failed flags out of
\[
1920;
\]

\item equivalently,
\[
32
\]
correctly closing flags;

\item no complete global \(r_3\);

\item no proof that the completed search space is infeasible.

\end{itemize}

The exact status is therefore

\[
\text{deeper bounded partial closure}
\]

but not

\[
\text{global transport theorem}.
\]

\subsection{Sigma boundary}

The unresolved inter-face sigma is not identified with:

\begin{itemize}

\item the signed \(G_{30}\to G_{15}\) cocycle;

\item the twelve-section binary closure class;

\item either integer parent;

\item the primitive registered-history cycle \(r_1\);

\item Thalean time;

\item or the normalized quadratic \(\widehat Q\).

\end{itemize}

An explicit structural crosswalk is required before any such
identification.

\subsection{Compatibility versus saturation}

The following distinction is locked:

\[
\text{local inadmissibility}
\neq
\text{global incompatibility}
\neq
\text{relational saturation}.
\]

In particular,

\[
\text{failure to find global sigma}
\]

is not evidence for

\[
\text{saturation}.
\]

The saturation conjecture applies only after a globally compatible
transport structure has been admitted.

\subsection{Finite summation ledger}

Paper 14 uses three established forms of finite accumulation.

\begin{center}
\begin{tabular}{p{0.27\linewidth}p{0.32\linewidth}p{0.31\linewidth}}
\toprule
Structure & Law & Domain \\
\midrule

incidence sum
&
\(
x\mapsto MM^{\mathsf T}x
\)
&
fifteen-sector real register
\\

binary regional sum
&
\(
d\delta=\omega
\)
&
twelve-section binary closure complex
\\

integer regional sum
&
\(
d\eta=\Omega
\)
&
integer refinement of closure complex
\\

history winding
&
\(
\tau_{\mathrm T}(\gamma_2\circ\gamma_1)
=
\tau_{\mathrm T}(\gamma_2)
+
\tau_{\mathrm T}(\gamma_1)
\)
&
registered-history cycle complex
\\

\bottomrule
\end{tabular}
\end{center}

These are related examples of lawful finite accumulation.

They are not one universal sum.

\subsection{Clean-release status}

The clean-release sequence

\[
\text{organization}
\to
\text{clean release}
\to
\text{overcompression}
\to
\text{contrast saturation}
\]

is an internal hypothesis.

The quadratic theorem proves only the contrast-concentrating direction.

No theorem in Paper 14 proves:

\begin{itemize}

\item a unique clean-release optimum;

\item a native loading parameter;

\item a universal release statistic;

\item or a physical emission mechanism.

\end{itemize}

\subsection{Empirical ledger}

The existing WMAP \(C_3/\Xi\) experiment is retained with its frozen
decision.

The exact empirical status is:

\[
\text{all four primary cells in predicted upper tail},
\]

\[
\text{three of four below }p=0.05,
\]

\[
p_{\mathrm{W,Gaussian}}
=
0.0539461,
\]

and therefore

\[
\boxed{\text{NO-GO REPLICATION}}.
\]

The experiment is not reinterpreted as a relational-saturation test.

Any saturation experiment requires a new frozen protocol.

\subsection{Physical bridge status}

The following remain physical bridge hypotheses rather than theorem
consequences:

\begin{itemize}

\item finite physical independent-distinction density;

\item minimum four-support of order \(\ell_P^4\);

\item physical gravity;

\item black-hole relational saturation;

\item horizon-like outward registration loss;

\item singularity avoidance;

\item black-hole entropy;

\item holography;

\item Hawking radiation;

\item a physical CMB saturation mechanism;

\item physical proper time;

\item Standard Model correspondence.

\end{itemize}

\subsection{Explicitly rejected identifications}

Paper 14 explicitly rejects the unsupported implications

\[
Q
=
\text{native temporal evolution},
\]

\[
\widehat Q
=
\text{physical dynamics},
\]

\[
\text{common mode}
=
\text{physical mass},
\]

\[
\text{centered magnitude}
=
\text{physical information},
\]

\[
\text{contrast concentration}
=
\text{gravity},
\]

\[
\text{projection-relative saturation}
=
\text{event horizon},
\]

\[
\text{finite saturation}
=
\text{singularity resolution},
\]

and

\[
\Xi
=
\text{saturation observable}.
\]

\subsection{Current theorem frontier}

The final status of Paper 14 is:

\[
\text{quadratic incidence law}
\quad
\text{proved},
\]

\[
\text{contrast concentration}
\quad
\text{proved},
\]

\[
\text{finite-stage non-erasure}
\quad
\text{proved},
\]

\[
\text{projection-relative saturation}
\quad
\text{defined},
\]

\[
\text{native loading crosswalk}
\quad
\text{open},
\]

\[
\text{global inter-face sigma}
\quad
\text{open},
\]

\[
\text{native relational-saturation theorem}
\quad
\text{open},
\]

\[
\text{physical saturation bridge}
\quad
\text{open}.
\]

\subsection{Final keeper}

The theorem frontier can be stated in one sentence:

\begin{quote}
Paper 14 proves how a finite linear register can lose relative
individuation without finite-stage state merger; the remaining problem
is to prove that an admitted native Thalean transport actually enters
that regime while retained complete history remains nontrivial.
\end{quote}


\bibliographystyle{plainnat}
\bibliography{refs}

\end{document}
